\documentclass[11pt,a4paper]{report}
% Adaptación de TiCSw/Uniandes-thesis-template-es.latex (master, 9a7f351).
% Se conserva el diseño de página y capítulos; citas y referencias en APA.
\usepackage[spanish,es-noshorthands,es-nodecimaldot]{babel}
\usepackage[T1]{fontenc}
\usepackage{lmodern,microtype}
\usepackage{graphicx}
\usepackage{fancyhdr,tocloft,vmargin}
\usepackage[font=small,format=plain,labelfont=bf]{caption}
\usepackage[centersections]{uniandestitlepage}
\usepackage{array,booktabs,longtable,tabularx,multirow}
\usepackage{amsmath,amssymb,setspace,enumitem,needspace,etoolbox}
\usepackage[authoryear,round]{natbib}
\usepackage{tikz}
\usetikzlibrary{arrows.meta,positioning,fit,calc}
\usepackage{xurl}
\usepackage[plainpages=false]{hyperref}

%\usepackage[papersize={17cm,24cm},inner=17mm,outer=17mm,top=20mm,bottom=20mm,includehead]{geometry}

\setmarginsrb  { 1.5in}  % left margin
                        { 0.8in}  % top margin
                        { 1.5in}  % right margin
                        { 0.8in}  % bottom margin
                        {  28pt}  % head height
                        {0.25in}  % head sep
                        {   9pt}  % foot height
                        { 0.3in}  % foot sep

% CONFIGURE CHAPTER TITLES
\usepackage[palatino]{quotchap}

% \usepackage[Lenny]{fncychap}
% \ChNameVar{\fontsize{14}{16}\usefont{OT1}{phv}{m}{n}\selectfont} 
% \ChNumVar{\fontsize{60}{62}\usefont{OT1}{ptm}{m}{n}\selectfont} 
% \ChTitleVar{\Huge\bfseries\rm}
% \ChRuleWidth{1pt}

% CONFIGURE SECTION TITLES
\usepackage[calcwidth]{titlesec}
\titleformat
{\section}[hang]{\sffamily\bfseries}
 {\Large\thesection}{12pt}{\Large}
 

% CONFIGURE PAGE HEADERS
\fancyhead{} % clear all header fields 
\fancyhead[R]{\thepage} 
\fancyhead[L]{\headerChapterTitle} 
 
\fancyfoot{} % clear all footer fields 
\fancyfoot[C]{} 

\renewcommand{\headrulewidth}{0.8pt} 

% Format the table of contents
\renewcommand{\contentsname}{Contenido} 
\renewcommand{\cfttoctitlefont}{ \hfill \huge}
\renewcommand{\cftloftitlefont}{ \hfill \huge }

\renewcommand{\cftlottitlefont}{ \hfill \huge }

\newcommand{\shortChapterTitle}{\leftmark}
\newcommand{\headerSectionTitle}{\parbox[b]{.86\textwidth}{\raggedleft\small\nouppercase{\rightmark}}}
\newcommand{\headerChapterTitle}{\parbox[b]{.86\textwidth}{\raggedright\small\nouppercase{\shortChapterTitle}}}

% CONFIGURE FIGURE CAPTIONS
\DeclareCaptionFont{ssblack}{\sffamily}
\DeclareCaptionFormat{figure}{\textbf{#1}~~#3}
\captionsetup[figure]{format=figure,labelfont=ssblack, textfont=ssblack, singlelinecheck=false, justification=centering}

% CONFIGURE TABLE CAPTIONS
\DeclareCaptionFormat{table}{\textbf{#1}~~#3}
\captionsetup[table]{format=table,labelfont=ssblack, textfont=ssblack, singlelinecheck=false, justification=centering}

\hypersetup{colorlinks=true,linkcolor=black,citecolor=black,urlcolor=black,pdftitle={Evaluación del movimiento de datos en flujos científicos con Pegasus bajo CondorIO},pdfauthor={Luis Sebastián Contreras Díaz}}
\urlstyle{same}
\AtBeginDocument{\renewcommand{\tablename}{Tabla}\renewcommand{\listtablename}{Índice de tablas}\renewcommand{\contentsname}{Contenido}}
\clubpenalty=10000\widowpenalty=10000\displaywidowpenalty=10000
\setlength{\emergencystretch}{2em}
\setlist{itemsep=2pt,topsep=4pt,leftmargin=*}
\newcolumntype{Y}{>{\raggedright\arraybackslash}X}
\newcolumntype{P}[1]{>{\raggedright\arraybackslash}p{#1}}
\newcommand{\MiB}{\,\mathrm{MiB}}
\newcommand{\DME}{\mathrm{DME}}
\newcommand{\Analyzer}{\textit{Analyzer}}
\newsavebox{\dmediagrambox}
\newcommand{\fitdmediagram}{\ifdim\wd\dmediagrambox>\linewidth\resizebox{\linewidth}{!}{\usebox{\dmediagrambox}}\else\usebox{\dmediagrambox}\fi}
\setlength{\bibhang}{1.27cm}
\setlength{\bibsep}{4pt}
\renewcommand{\bibfont}{\normalfont\normalsize}
\renewcommand{\bibsection}{\chapter*{Referencias}\markboth{Referencias}{Referencias}\addcontentsline{toc}{chapter}{Referencias}}
\setcounter{tocdepth}{1}
\setcounter{secnumdepth}{2}
\AtBeginEnvironment{longtable}{\singlespacing}

% Reversos de cortesía sin encabezado ni numeración visibles.
\makeatletter
\renewcommand{\cleardoublepage}{\clearpage\if@twoside\ifodd\c@page\else\hbox{}\thispagestyle{empty}\newpage\if@twocolumn\hbox{}\newpage\fi\fi\fi}
\makeatother
\raggedbottom

% Alinear arriba las tablas que ocupen una página de flotantes.
\makeatletter\setlength{\@fptop}{0pt}\setlength{\@fpsep}{12pt}\makeatother

\begin{document}
\hypersetup{pageanchor=false}
% Datos reales de la tesis. Portada basada en other/title.tex de TiCSw.
\sectonesize{20pt}\sectoneskip{25pt}
\secttwosize{16pt}\secttwoskip{20pt}
\sectthreesize{11pt}\sectthreeskip{15pt}
\sectone{\centering Evaluación del movimiento de datos en flujos científicos con Pegasus bajo CondorIO\par\vspace{5mm}{\fontsize{12pt}{16pt}\selectfont Definición, implementación y evaluación de DME bajo CondorIO\par}}
\secttwo{\centering Luis Sebastián Contreras Díaz\par\vspace{2mm}{\fontsize{11pt}{14pt}\selectfont Código 202311819\par}}
\sectthree{\centering Trabajo de grado para optar al título de\\Ingeniero de Sistemas y Computación\par\vspace{7mm}\textbf{Asesores}\par Aurelio Antonio Vivas Meza, Ph. D.\\Argonne National Laboratory\par\vspace{3mm}Harold Enrique Castro Barrera, Ph. D.\\Universidad de los Andes\par}
\sectfour{}
\date{\centering Bogotá, Colombia\\2026}
\maketitle

\clearpage\pagenumbering{roman}
\hypersetup{pageanchor=true}
\chapter*{Resumen}
La eficiencia del movimiento de datos (DME, por sus siglas en inglés) relaciona el volumen que requieren las dependencias de un flujo científico con sus entregas completas corroboradas. La referencia depende de la ubicación efectiva y del relevo obligatorio bajo CondorIO. La evidencia histórica detalla 45 ejecuciones de quince condiciones en cinco patrones, más nueve comparaciones de ubicación libre y un piloto. Una campaña prospectiva completa 63 ejecuciones principales en doce condiciones, separadas de doce pilotos válidos y un intento técnico fallido preservado. En seis repeticiones por condición sintética, agrupar reduce el observado de 60 a 40 MiB en el pipeline y de 100 a 70 MiB con un archivo compartido; las particiones distintas conservan 40 MiB. Una intervención sobre el retorno de intermedios lleva el pipeline a 20 MiB y DME igual a uno sin alterar el final. Tres repeticiones por variante de un subworkflow de Montage confirman agrupación y retorno selectivo, con finales científicamente equivalentes. DME normaliza respecto de las obligaciones; con requerido común conserva el orden inverso del observado y entre ubicaciones no ordena bytes por sí sola. Un contraejemplo analítico muestra que el exceso no siempre es evitable bajo jobs aislados sin caché. El costo de captura nativa sigue sin aislarse.

\smallskip\noindent\textit{Palabras clave.} Flujos de trabajo científicos, movimiento de datos, Pegasus, HTCondor, agrupación de tareas.

\section*{Abstract}
Data Movement Efficiency (DME) relates required workflow volume to corroborated complete deliveries. Its reference depends on the graph, placement, and mandatory CondorIO relay. Historical validation covers 45 runs across fifteen conditions and five patterns, plus nine free-placement comparisons and one pilot. A prospective campaign completed 63 main runs across twelve conditions, with twelve valid pilots and one preserved technical failure. Across six repetitions per synthetic condition, clustering reduced observed volume from 60 to 40 MiB in the pipeline and from 100 to 70 MiB for a shared file; distinct partitions retained 40 MiB. Omitting intermediate returns further reduced the pipeline to 20 MiB and DME equal to one, without changing its final output. Three runs per Montage variant confirmed these mechanisms and equivalent outputs. With common required volume, DME preserves the inverse ordering of observed bytes; across placements, it does not rank bytes alone. An analytical counterexample shows an unattainable bound for one graph with isolated, no-cache jobs. Native capture overhead remains unisolated.

\smallskip\noindent\textit{Keywords.} Scientific workflows, data movement, Pegasus, HTCondor, task clustering.
\begingroup
\setlength{\cftbeforechapskip}{6pt}
\tableofcontents
\endgroup
\clearpage\pagenumbering{arabic}\pagestyle{fancy}
\chapter{Introducción}
\section{Motivación}
Un flujo de trabajo científico coordina tareas que consumen y producen archivos. Pegasus prepara el plan y HTCondor ejecuta sus unidades de trabajo, llamadas \textit{jobs}. El coordinador gestiona el flujo y recibe o envía sus archivos; W1 y W2 designan los dos nodos de ejecución (workers). En una cadena de tareas (\textit{pipeline}), cada etapa puede ejecutarse en un job separado o compartirlo con otras etapas. Esa decisión puede cambiar cuántas veces se entrega un archivo intermedio.

Esta tesis propone una forma de medir ese cambio y explicar qué entregas lo producen. Cuando el usuario agrupa dos etapas, puede dejar de reenviarse un intermedio al segundo job y conservarse su retorno al coordinador. Para distinguir esas operaciones hay que cotejar las dependencias, la ubicación de las tareas, el tamaño de los archivos y los registros de ejecución. En un caso pequeño se puede hacer la cuenta a mano. Aplicar el mismo criterio a varias configuraciones requiere un procedimiento explícito.

\section{Problema}
Un total de bytes carece de contexto si se desconoce qué archivos debían estar disponibles y en qué máquinas. Los registros de transferencia también pueden incluir ejecutables y otros archivos auxiliares. Para medir el volumen científico se necesita identificar los archivos que pertenecen al flujo y distinguirlos de los que requiere el sistema para ejecutarlo.

Agrupar tareas puede cambiar las entregas, y cambiar su ubicación también puede modificar las necesidades del flujo. La comparación debe separar esos dos efectos. Si un intermedio cruza entre workers, CondorIO exige su relevo por el coordinador. Si productor y consumidores permanecen en el mismo worker, sus contratos aún pueden declarar retornos y reenvíos que exceden esa cota.

\Needspace{5\baselineskip}
\textbf{Pregunta de investigación.} Dado un flujo de trabajo y la asignación real de sus tareas a máquinas, \textquestiondown{}qué cota inferior de movimiento de archivos imponen sus dependencias y el relevo obligatorio por el coordinador bajo CondorIO, y cómo se relaciona con las entregas completas corroboradas?

\section{Justificación}
La eficiencia del movimiento de datos (DME, por sus siglas en inglés) relaciona el volumen observado con las obligaciones que impone el flujo en su ubicación efectiva. Con las mismas necesidades, una reducción del observado indica que se entregaron menos bytes. Si cambian las necesidades por una nueva ubicación, la razón puede aumentar aunque el volumen entregado siga igual. Conservar las fuentes de ambos cálculos permite distinguir estas situaciones.

El aporte es una referencia por archivo y destino que permite interpretar el volumen entregado. La agrupación sirve para contrastarla: puede eliminar reenvíos, permitir compartir una entrada o conservar las mismas entregas. La evaluación distingue esos mecanismos y sus límites. Analyzer implementa el cálculo y vincula sus resultados con las fuentes que los sustentan.

\section{Objetivos}
\textbf{Objetivo general.} Diseñar, implementar y validar una herramienta que permita calcular una métrica para evaluar la eficiencia del movimiento de datos durante la ejecución de flujos de trabajo gestionados por Pegasus, utilizando los patrones process, pipeline, data aggregation, data distribution y data redistribution para evaluar la sensibilidad de la métrica bajo condiciones controladas de ejecución.

\textbf{Objetivos específicos.}
\begin{enumerate}
\item Caracterizar, a partir de la literatura y de la documentación técnica de Pegasus y HTCondor, las variables y fuentes de información necesarias para medir el movimiento de datos y la estructura de los patrones process, pipeline, data aggregation, data distribution y data redistribution.
\item Definir una métrica para evaluar la eficiencia del movimiento de datos, especificando sus variables, unidades, procedimiento de cálculo, supuestos, condiciones de aplicación y comportamiento esperado ante los patrones seleccionados.
\item Diseñar e implementar una herramienta que automatice la extracción, integración y procesamiento de la información requerida para calcular la métrica definida a partir de ejecuciones gestionadas por Pegasus.
\item Validar la métrica y la herramienta mediante ejecuciones controladas de los cinco patrones de movimiento de datos, documentando la corrección y consistencia del cálculo, la sensibilidad de la métrica, las condiciones de aplicación, el sobrecosto de instrumentación, las fortalezas y las limitaciones de la solución.
\end{enumerate}


\Needspace{5\baselineskip}
\section{Alcance}
La tesis abarca la especificación de DME, su cálculo con \Analyzer{} y su evaluación en flujos gestionados por Pegasus y HTCondor con CondorIO. El modelo distingue la cota que exigen el grafo y el relevo por el coordinador con una asignación dada del volumen de las entregas corroboradas. La herramienta integra el plan, el historial y los registros de transferencia, presenta el desglose por archivo y destino, y señala cuándo la evidencia es insuficiente.

La medición comprende las entradas científicas, los intermedios y los resultados declarados, desde las entradas disponibles en el coordinador hasta la recuperación de los resultados. La unidad de ubicación es la máquina virtual y los archivos deben conservar su identidad y tamaño. Se contabilizan entregas completas. Los registros de archivos auxiliares se usan para conciliar los contadores de transferencia. Los volúmenes se expresan en bytes y DME es su cociente adimensional. Para aplicar el procedimiento a otro flujo se requiere evidencia de sus archivos, ubicaciones y entregas bajo el mismo perfil de medición.

El bloque A reúne 24 condiciones y 72 ejecuciones elegibles. La validación histórica detallada analiza 45 ejecuciones de 15 condiciones en los cinco patrones. En un contraste separado se comparan L01, GL2 y GL3 con ubicación libre (n=3) y se documenta un piloto GL2; las asignaciones efectivas se recuperan del historial. Se conservan también cuatro ejecuciones históricas de Montage con ubicación W1/W2. Una campaña prospectiva separada completa 63 ejecuciones principales: seis por condición sintética y tres por variante de Montage con ambas tareas en W1. Doce pilotos válidos y un intento técnico fallido se documentan aparte. Se comparan agrupaciones y una intervención adicional en el retorno de intermedios, conservando los finales científicos. La duración es una respuesta secundaria. El costo de generar registros durante la ejecución queda por aislar.

\section{Contenido y relación con los objetivos}
El capítulo~\ref{ch:marco} presenta los patrones, conceptos y fuentes del primer objetivo. El capítulo~\ref{ch:diseno} define la métrica y sus condiciones de aplicación, correspondientes al segundo. El capítulo~\ref{ch:implementacion} documenta la herramienta del tercero. El capítulo~\ref{ch:validacion} aborda el cuarto mediante la comprobación del cálculo, la consistencia, el contraste libre y la campaña prospectiva de agrupación. Los capítulos finales presentan las conclusiones, el trabajo futuro y los materiales de reproducción.

\begin{table}[htbp]\centering\small
\caption{Evidencia asociada a los objetivos específicos}
\label{tab:objetivos_evidencia}
\begin{tabularx}{\linewidth}{@{}P{2.5cm}YY@{}}
\toprule Objetivo & Evidencia en la tesis & Qué se puede comprobar\\\midrule
OE1. Caracterización & Definiciones, cinco grafos y fuentes en el marco teórico y la implementación. & Qué cambia entre tareas, jobs, ubicaciones y archivos, y de dónde se obtiene cada variable.\\
OE2. Definición & Cota por archivo y destino con relevo, unidades, cobertura y estados. & Cuándo se puede calcular DME y cuándo el observado o la referencia explican el cambio.\\
OE3. Implementación & Arquitectura, contratos, diagnósticos y materiales de reproducción. & Cómo se obtienen desgloses verificables. La regla de relevo está implementada.\\
OE4. Validación & Cotejo histórico de 45 ejecuciones, nueve libres y un piloto; 63 principales prospectivas en doce condiciones, con pilotos separados y control de finales. & Correctitud, consistencia y sensibilidad bajo el perfil ensayado. Queda por aislar el costo de captura nativa.\\
\bottomrule
\end{tabularx}
\end{table}

\chapter{Marco teórico}
\label{ch:marco}
\section{Tareas, archivos y jobs}
Un flujo de trabajo científico puede representarse como un grafo dirigido acíclico. Cada tarea produce archivos que consumen las tareas sucesoras. Las entradas provienen de fuera del grafo; los resultados finales deben quedar disponibles en el destino declarado. Los archivos que conectan tareas son intermedios.

Pegasus transforma el flujo científico en un plan ejecutable con jobs científicos y operaciones auxiliares \citep{deelman2015}. Cada job contiene una tarea o un grupo. DAGMan (\textit{Directed Acyclic Graph Manager}) respeta las dependencias entre jobs \citep{htcondordagman}. HTCondor asigna cada job a recursos elegibles y registra su ejecución \citep{htcondorintro}. El plan muestra qué tareas se agruparon; el historial, dónde se ejecutaron.

\section{Ubicación y entrega de archivos}
En CondorIO, el coordinador prepara las entradas y recibe las salidas. Cada job trabaja en un directorio aislado o \textit{sandbox}. Un archivo puede volver al coordinador y salir hacia otro job, incluso si ambos corren en la misma máquina \citep{htcondorordinary,pegasustransfers}. Al agrupar tareas consecutivas, algunas lecturas pueden resolverse en el mismo sandbox. Como las salidas declaradas pueden conservarse, los registros permiten comprobar el ahorro.

El volumen observado suma el tamaño de cada entrega completa; un mismo archivo puede generar varias entregas. Los contadores nativos también pueden incluir ejecutables, metadatos y registros auxiliares. Por eso se conservan los archivos y tamaños científicos, el contrato efectivo de cada job, sus estadísticas y eventos. Kickstart y los metadatos de Pegasus aportan información de procedencia \citep{pegasusprovenance}.

\section{Patrones de flujo}
Los patrones representan relaciones de producción y consumo, y permiten aislar cambios en el movimiento de archivos. La caracterización de workflows científicos identifica estructuras y dependencias diversas \citep{bharathi2008}. Aquí, procesamiento sirve como control; pipeline sigue dependencias consecutivas; agregación reúne salidas; distribución divide un conjunto; y redistribución combina reunión y reparto.
En las Figuras~\ref{fig:pat-process} a~\ref{fig:pat-redistribution}, los círculos son tareas y los rectángulos archivos; las flechas indican producción y consumo, no transferencias entre máquinas. La ubicación y los contratos se incorporan después al cálculo. Las condiciones comparadas se describen en la sección~\ref{sec:repro_patrones}.

\subsection{Procesamiento}
Una tarea transforma una entrada en un resultado final. La Figura~\ref{fig:pat-process} presenta el caso utilizado en P01 y P02.

\par\medskip\noindent\begin{minipage}{\linewidth}
\captionsetup{type=figure}
\caption{Patrón de procesamiento}
\label{fig:pat-process}
\centering
\begin{lrbox}{\dmediagrambox}\begin{tikzpicture}[>=Stealth,font=\small,
 t/.style={draw,circle,minimum size=0.75cm,inner sep=1pt},
 f/.style={draw,rectangle,minimum width=1cm,minimum height=0.6cm,inner sep=2pt}]
\node[f] (a) at (0,0) {$A$};
\node[t] (t) at (3,0) {$T_1$};
\node[f] (b) at (6,0) {$B$};
\draw[->] (a)--(t); \draw[->] (t)--(b);
\end{tikzpicture}\end{lrbox}\fitdmediagram
\par\smallskip\raggedright\normalsize\textit{Nota.} $A$ es la entrada externa y $B$ el resultado final. Ambos miden 10 MiB en P01 y 20 MiB en P02. La estructura permanece constante al cambiar el tamaño.
\end{minipage}\par\medskip

\subsection{Pipeline}
Cada etapa consume la salida de la anterior. La Figura~\ref{fig:pat-pipeline} conserva la identidad de los intermediarios, aunque las tareas se agrupen después en jobs.

\par\medskip\noindent\begin{minipage}{\linewidth}
\captionsetup{type=figure}
\caption{Patrón pipeline y sus archivos intermedios}
\label{fig:pat-pipeline}
\centering
\begin{lrbox}{\dmediagrambox}\begin{tikzpicture}[>=Stealth,font=\small,
 t/.style={draw,circle,minimum size=0.75cm,inner sep=1pt},
 f/.style={draw,rectangle,minimum width=1cm,minimum height=0.6cm,inner sep=2pt}]
\node[f] (a) at (0,0) {$A$};
\node[t] (t1) at (2,0) {$T_1$};
\node[f] (f12) at (4,0) {$f_{12}$};
\node[t] (t2) at (6,0) {$T_2$};
\node[f] (f23) at (8,0) {$f_{23}$};
\node[t] (t3) at (10,0) {$T_3$};
\node[f] (b) at (12,0) {$B$};
\foreach \u/\v in {a/t1,t1/f12,f12/t2,t2/f23,f23/t3,t3/b}{\draw[->] (\u)--(\v);}
\end{tikzpicture}\end{lrbox}\fitdmediagram
\par\smallskip\raggedright\normalsize\textit{Nota.} En L01, $A$, $f_{12}$, $f_{23}$ y $B$ miden 10 MiB. Agrupar tareas conserva estos archivos y sus dependencias, aunque puede modificar las entregas del mecanismo de ejecución.
\end{minipage}\par\medskip

\Needspace{16\baselineskip}
\subsection{Agregación}
Una tarea agregadora consume archivos distintos de varios productores. La Figura~\ref{fig:pat-aggregation} muestra los cuatro productores del caso base.

\par\medskip\noindent\begin{minipage}{\linewidth}
\captionsetup{type=figure}
\caption{Patrón de agregación: varios productores alimentan una tarea agregadora}
\label{fig:pat-aggregation}
\centering
\begin{lrbox}{\dmediagrambox}\begin{tikzpicture}[>=Stealth,font=\small,
 t/.style={draw,circle,minimum size=0.75cm,inner sep=1pt},
 f/.style={draw,rectangle,minimum width=1cm,minimum height=0.6cm,inner sep=2pt}]
\node[t] (agg) at (7,0) {$G$};
\node[f] (b) at (10,0) {$B$};
\foreach \i/\y in {1/1.8,2/0.6,3/-0.6,4/-1.8}{
 \node[f] (a\i) at (0,\y) {$A_{\i}$};
 \node[t] (p\i) at (2,\y) {$P_{\i}$};
 \node[f] (f\i) at (4,\y) {$f_{\i}$};
 \draw[->] (a\i)--(p\i); \draw[->] (p\i)--(f\i); \draw[->] (f\i)--(agg);
}
\draw[->] (agg)--(b);
\end{tikzpicture}\end{lrbox}\fitdmediagram
\par\smallskip\raggedright\normalsize\textit{Nota.} En A01, cada entrada $A_i$ y cada intermediario $f_i$ mide 10 MiB; $G$ produce el final $B$ de 40 MiB. Cada subíndice identifica un archivo distinto.
\end{minipage}\par\medskip

\subsection{Distribución}
Un productor puede repartir particiones distintas o compartir un archivo entre varios consumidores. La Figura~\ref{fig:pat-distribution} distingue ambas relaciones, necesarias para interpretar el efecto de agrupar.

\par\medskip\noindent\begin{minipage}{\linewidth}
\captionsetup{type=figure}
\caption{Patrón de distribución: particiones distintas y archivo compartido}
\label{fig:pat-distribution}
\centering
\begin{lrbox}{\dmediagrambox}\begin{tikzpicture}[>=Stealth,font=\small,
 t/.style={draw,circle,minimum size=0.7cm,inner sep=1pt},
 f/.style={draw,rectangle,minimum width=0.95cm,minimum height=0.55cm,inner sep=2pt}]
\node[anchor=west] at (-0.5,2.5) {(a) Particiones distintas};
\node[f] (a) at (0,0) {$A$}; \node[t] (s) at (2,0) {$S$};
\draw[->] (a)--(s);
\foreach \i/\y in {1/1.65,2/0.55,3/-0.55,4/-1.65}{
 \node[f] (q\i) at (5,\y) {$q_{\i}$};
 \node[t] (c\i) at (8,\y) {$C_{\i}$};
 \node[f] (b\i) at (11,\y) {$B_{\i}$};
 \draw[->] (s)--(q\i); \draw[->] (q\i)--(c\i); \draw[->] (c\i)--(b\i);
}
\node[anchor=west] at (-0.5,-3) {(b) Archivo compartido};
\node[f] (aa) at (0,-5.5) {$A$}; \node[t] (p) at (2,-5.5) {$P$};
\node[f] (g) at (5,-5.5) {$g$};
\draw[->] (aa)--(p); \draw[->] (p)--(g);
\foreach \i/\y in {1/-3.85,2/-4.95,3/-6.05,4/-7.15}{
 \node[t] (cc\i) at (8,\y) {$C_{\i}$};
 \node[f] (bb\i) at (11,\y) {$B_{\i}$};
 \draw[->] (g)--(cc\i); \draw[->] (cc\i)--(bb\i);
}
\end{tikzpicture}\end{lrbox}\fitdmediagram
\par\smallskip\raggedright\normalsize\textit{Nota.} En D02 y GD4, $S$ divide una entrada de 10 MiB en cuatro particiones $q_i$ de 2,5 MiB y cada final $B_i$ conserva ese tamaño. En AC1, AC2 y AC4, $P$ produce un único archivo $g$ de 10 MiB, consumido por cuatro tareas que generan finales de 10 MiB. Las cuatro aristas de consumo de $g$ conservan una sola identidad de archivo.
\end{minipage}\par\medskip

\subsection{Redistribución}
La redistribución agrega archivos de varios productores y divide el resultado para nuevos consumidores. La Figura~\ref{fig:pat-redistribution} representa ambas operaciones y sus intermediarios.

\par\medskip\noindent\begin{minipage}{\linewidth}
\captionsetup{type=figure}
\caption{Patrón de redistribución: agregación seguida de partición y consumo}
\label{fig:pat-redistribution}
\centering
\begin{lrbox}{\dmediagrambox}\begin{tikzpicture}[>=Stealth,font=\small,
 t/.style={draw,circle,minimum size=0.65cm,inner sep=1pt},
 f/.style={draw,rectangle,minimum width=0.85cm,minimum height=0.55cm,inner sep=2pt}]
\node[t] (agg) at (5.25,0) {$G$};
\node[f] (h) at (7,0) {$h$};
\node[t] (s) at (8.75,0) {$S$};
\draw[->] (agg)--(h); \draw[->] (h)--(s);
\foreach \i/\y in {1/1.8,2/0.6,3/-0.6,4/-1.8}{
 \node[f] (a\i) at (0,\y) {$A_{\i}$};
 \node[t] (p\i) at (1.75,\y) {$P_{\i}$};
 \node[f] (f\i) at (3.5,\y) {$f_{\i}$};
 \draw[->] (a\i)--(p\i); \draw[->] (p\i)--(f\i); \draw[->] (f\i)--(agg);
 \node[f] (q\i) at (10.5,\y) {$q_{\i}$};
 \node[t] (c\i) at (12.25,\y) {$C_{\i}$};
 \node[f] (b\i) at (14,\y) {$B_{\i}$};
 \draw[->] (s)--(q\i); \draw[->] (q\i)--(c\i); \draw[->] (c\i)--(b\i);
}
\end{tikzpicture}\end{lrbox}\fitdmediagram
\par\smallskip\raggedright\normalsize\textit{Nota.} En R01-E, las entradas $A_i$, los intermediarios $f_i$, las particiones $q_i$ y los finales $B_i$ miden 10 MiB cada uno. El archivo agregado $h$ mide 40 MiB. Las tareas $G$ y $S$ representan agregación y división, respectivamente.
\end{minipage}\par\medskip


\section{Antecedentes y relación con DME}
\citet{chen2015} estudian la agrupación de tareas y sus efectos en la ejecución, una decisión que modifica el plan. Esta tesis pregunta qué entregas científicas cambian entre configuraciones y cómo se relacionan con las obligaciones del mismo grafo. Para responderla, se requiere el mapa de tareas a jobs y las listas de archivos transferidos.

Pietri y Sakellariou (\citeyear{pietri2018}) estudian la planificación para reducir comunicación. \citet{fortmeier2013} relacionan modelos de volumen con métricas de partición de hipergrafos. Aquí, $n_f$ cuenta, para cada archivo, los destinos distintos del origen y se extiende con el relevo CondorIO. No se adopta su métrica de comunicación entre todos los vecinos: se fija la ubicación observada y no se optimiza la partición.

La métrica MUE (Memory Usage Efficiency) usa una referencia teórica de accesos a memoria y el aprovechamiento del ancho de banda para evaluar una implementación \citep{fuhrer2018}. Este antecedente ilustra cómo una referencia ayuda a interpretar un volumen. DME opera a otra escala: define necesidades por archivo y ubicación y las contrasta con las entregas corroboradas durante una ejecución del flujo.

\texttt{pegasus-statistics} informa conteos y tiempos de tareas y jobs \citep{pegasusstats}, pero no expresa la cota por archivo, ubicación y relevo definida aquí. Analyzer añade esa referencia y concilia las entregas científicas. La integración de registros de distintas capas tiene antecedentes en \citet{papadimitriou2021}; para compararlos, las fuentes deben corresponder a una misma ejecución.

\begin{table}[htbp]\centering\small
\caption{Variables y fuentes de la caracterización}
\begin{tabularx}{\linewidth}{@{}P{2.75cm}YY@{}}
\toprule Variable & Fuente de referencia & Función en la medición\\\midrule
Rol del archivo & Grafo científico y declaración de finales. & Distinguir entrada, intermedio y final, y determinar productores y consumidores.\\
Identidad y tamaño & Catálogos y metadatos de ejecución. & Vincular cada archivo y expresarlo en bytes.\\
Tarea y job & Plan de Pegasus y descripción del grupo. & Identificar si dos tareas usan el mismo espacio de ejecución.\\
Máquina e intento & Historial y eventos de HTCondor. & Resolver ubicación efectiva, terminación y alcance del intento.\\
Contrato de entrega & Listas efectivas de entrada y salida. & Identificar qué archivos pueden contabilizarse en cada dirección.\\
Contadores y tiempos & Estadísticas de transferencia y métricas de DAGMan. & Conciliar bytes y conteos. Medir duración del workflow.\\
\bottomrule
\end{tabularx}
\end{table}

Distinguir el archivo lógico de cada copia entregada evita contar solo por nombre de tarea o job. En el pipeline, un intermedio puede ser local según la referencia si productor y consumidor están en W1. Si sus jobs usan sandboxes distintos, el mecanismo puede devolverlo al coordinador y entregarlo al siguiente job. El plan y los registros permiten establecer ese recorrido; el grafo y la ubicación, determinar si el dato hacía falta en otra máquina.

\chapter{Análisis y diseño de la solución}
\label{ch:diseno}
\section{Especificación y frontera}
La solución debe indicar qué cota de volumen imponen las dependencias y el relevo bajo CondorIO, cuánto se entregó y qué archivos y destinos explican la diferencia. Los tres resultados se refieren al mismo grafo y a la misma ejecución. El registro de salida vincula cada obligación o entrega con el archivo y el contexto utilizado para contarlo.

La medición empieza con las entradas científicas en el coordinador y termina con la disponibilidad de los finales declarados. El observado suma el tamaño lógico de cada archivo científico por entrega completa corroborada. Los auxiliares sirven para conciliar registros y se excluyen de ese volumen. DME no mide cabeceras, retransmisiones ni el efecto de comprimir el transporte. Una ubicación corresponde a una máquina virtual. El perfil de medición es el conjunto de condiciones bajo las cuales se reconstruye el cálculo. Requiere identidades y tamaños estables, un productor por archivo generado, un intento efectivo por job y evidencia suficiente de los contratos. Ante fuentes incompatibles o incompletas se conserva el diagnóstico.

\section{Modelo de tareas, archivos y ubicación}
Se representa el flujo científico mediante un grafo dirigido acíclico $G=(T,F,D)$, donde $T$ es el conjunto de tareas, $F$ el de archivos y $D$ contiene aristas de producción $(t,f)$ y de consumo $(f,t)$. Cada archivo generado tiene un productor y puede tener varios consumidores. Las entradas externas se reúnen en $F_{\mathrm{ext}}$ y los resultados que deben quedar en el coordinador en $F_{\mathrm{fin}}$. Sea $m$ el coordinador, $p(f)$ el productor de un archivo generado y $\pi(t)$ la máquina donde se ejecutó la tarea $t$. El mapa tarea--job se obtiene del plan, mientras que la ubicación se comprueba en el historial del job.

Cada archivo tiene un tamaño $s(f)$ no negativo en bytes. Un MiB equivale a $2^{20}$ bytes. El origen $o(f)$ de una entrada externa es $m$. Para un archivo generado es $\pi(p(f))$. El conjunto de destinos requeridos es
\begin{equation}
R_f=\{\pi(t)\mid (f,t)\in D\}\cup Q_f.
\end{equation}
Se define $Q_f=\{m\}$ si $f$ es un final declarado o si es generado y algún consumidor se ejecuta en un worker distinto del productor; en otro caso, $Q_f=\varnothing$. Bajo CondorIO se consideran solo transferencias coordinador--worker: llevar un archivo de W1 a W2 exige recibirlo primero en el coordinador. Los finales y el relevo comparten el mismo destino $m$, que se cuenta una sola vez. Las entradas externas ya tienen origen en $m$.

\section{Definición de la métrica}
Para cada archivo $f$, sea $n_f=\lvert R_f\setminus\{o(f)\}\rvert$ el número de máquinas distintas de su origen en las que debe estar disponible. Cada máquina aporta una unidad aunque allí lo consuman varias tareas. Si todos los destinos están en el origen o $R_f$ es vacío, $n_f=0$. Así, $s(f)n_f$ expresa los bytes requeridos para el archivo $f$. Los dos volúmenes se calculan como
\begin{align}
M_{\mathrm{req}}&=\sum_{f\in F}s(f)\,n_f,\label{eq:requerido}\\
M_{\mathrm{obs}}&=\sum_{e\in\mathcal E_{\mathrm{conf}}}s(f_e),\label{eq:observado}
\end{align}
donde $\mathcal E_{\mathrm{conf}}$ reúne las entregas científicas completas corroboradas y $f_e$ es el archivo de la entrega $e$. La referencia es una cota inferior con ubicación ejecutada fija, archivos completos, transferencias solo por el coordinador y una entrega por worker de destino. La cota se alcanza en las ejecuciones con observado igual al requerido, entre ellas A02, D02, GD4, AC4 y R02-E. Su alcanzabilidad no está demostrada para grafos generales ni para los casos con exceso: no modela todas las restricciones de sandboxes y contratos.

\begin{table}[htbp]\centering\small
\caption{Conteo de destinos con relevo obligatorio}\label{tab:ejemplos_relevo}
\begin{tabularx}{\linewidth}{@{}Ylll@{}}
\toprule Archivo & $Q_f$ & $R_f\setminus\{o(f)\}$ & $n_f$\\\midrule
Nace en W1, se usa en W2 & $\{m\}$ & $\{m,\mathrm{W2}\}$ & 2\\
Nace y se usa en W1 & $\varnothing$ & $\varnothing$ & 0\\
Final en W1, consumido en W2 & $\{m\}$ & $\{m,\mathrm{W2}\}$ & 2\\
Final nacido en W1 & $\{m\}$ & $\{m\}$ & 1\\
Entrada en $m$, usada en W1 & $\varnothing$ & $\{\mathrm{W1}\}$ & 1\\
\bottomrule\end{tabularx}
\end{table}

La métrica resume la relación entre ambos volúmenes:
\begin{equation}
\DME=\frac{M_{\mathrm{req}}}{M_{\mathrm{obs}}}\quad (M_{\mathrm{obs}}>0).
\label{eq:dme}
\end{equation}
El cociente es adimensional. El análisis también usa la diferencia $M_{\mathrm{obs}}-M_{\mathrm{req}}$ cuando ayuda a explicar las entregas que separan la ejecución de su referencia. Esta diferencia expresa exceso sobre la cota inferior. DME igual a uno significa que el observado está en la cota; un valor menor indica exceso sobre ella. El exceso puede incluir retornos y reenvíos entre sandboxes. Agrupar puede reducirlos, según los contratos, pero no se ha demostrado que todo exceso sea evitable.

Un contraejemplo analítico ilustra este límite bajo restricciones adicionales. Sean las dependencias T1--T2--T3--T4 y T1--T4, con ubicación W1/W2/W1/W2. La entrada, los tres intermedios y el final tienen tamaño unitario; el primer intermedio alimenta T2 y T4. La referencia por archivo y destino, con relevo obligatorio por el coordinador, suma ocho unidades. Si los jobs son secuenciales y aislados, los archivos completos pasan por el coordinador y no se reutilizan entre jobs, T2 y T4 requieren dos envíos del intermedio: el mínimo observado es nueve unidades y DME=8/9. Las únicas fusiones no triviales (T1/T3 o T2/T4) contraen caminos y crean ciclos. El ejemplo es analítico, no ejecutado; depende de la ausencia de caché o reutilización y no establece un límite universal.

Con $M_{\mathrm{req}}$ positivo y común, una DME mayor equivale a un $M_{\mathrm{obs}}$ menor: no añade una clasificación distinta por bytes. Su aporte es normalizar el observado respecto de las obligaciones y, junto con el desglose, explicar qué entregas forman la diferencia. Por ello, DME menor que uno no demuestra que el exceso sea evitable.

\subsection{Cobertura y estados de cálculo}
Para probar la cota, se exige que cada obligación $(f,v)$ pueda asociarse a una entrega corroborada del mismo archivo al mismo destino. Dos obligaciones distintas deben asociarse a entregas distintas. En el perfil CondorIO, cada archivo generado que llega a otro worker pasa antes por el coordinador; por ello existen una entrega de relevo y otra al destino. Esta necesidad del mecanismo no sustituye su corroboración en los registros. Si se cumple esta condición, el requerido suma los tamaños de un subconjunto de las entregas observadas. Por tanto, $0\leq M_{\mathrm{req}}\leq M_{\mathrm{obs}}$, y $0\leq\DME\leq 1$ cuando el denominador es positivo. La cobertura se verifica por identidad y destino, aun con archivos de cero bytes. El observado completo exige incluir todas las entregas científicas del intento admitido y distinguir cada entrega de registros duplicados.

Si ambos volúmenes son cero, la razón queda indefinida. Con requerido cero y observado positivo, el valor es cero. Con requerido positivo y observado cero, debe revisarse la cobertura antes de interpretar la ejecución. Si falta evidencia para determinar un volumen, se conserva su estado indeterminado. Un cociente mayor que uno contradice la condición de cobertura y exige examinar la asignación, las obligaciones y las entregas incluidas.

\section{Interpretación de los contrastes}
La línea base experimental es la configuración de comparación. La cota con relevo se calcula para cada ejecución. Para un mismo trabajo y perfil, sean $O_b$ y $O_v$ los observados de la base y la variante. La reducción relativa es
\begin{equation}
\rho=\frac{O_b-O_v}{O_b}\qquad(O_b>0).
\end{equation}
El producto $100\rho$ expresa el porcentaje de reducción respecto de la base. Si el requerido es común y positivo, una reducción del observado aumenta DME. El desglose permite identificar las entregas que explican esa reducción.

Si también cambia el requerido, sean $R_b$ y $R_v$ sus valores respectivos. Con los cuatro volúmenes positivos, se puede escribir
\begin{equation}
\frac{\DME_v}{\DME_b}=\frac{R_v/R_b}{O_v/O_b}.
\end{equation}
La expresión separa ambos cambios relativos: DME puede aumentar sin ahorro si $M_{\mathrm{obs}}$ permanece constante y crece el requerido. Si cambia el requerido, se comparan los observados para ordenar bytes.



GL3 y GD4 observan 40 MiB con requeridos de 20 y 40 MiB: DME vale $1/2$ y $1$. La referencia distingue esa relación, sin establecer superioridad global entre trabajos distintos.

La duración se registra por separado: agrupar también cambia los jobs y la actividad administrativa.

\section{Arquitectura del cálculo}
La Figura~\ref{fig:arquitectura} distingue las fuentes que sustentan la referencia de las que corroboran las entregas. Las dos ramas comparten la identificación de la ejecución, los archivos y sus tamaños. Al integrarlas se comprueban consistencia y cobertura antes de emitir la razón y sus desgloses.

\begin{figure}[htbp]
\caption{Fuentes y cálculo de los dos volúmenes}
\label{fig:arquitectura}
\centering
\begin{lrbox}{\dmediagrambox}\begin{tikzpicture}[>=Stealth,font=\small]
\tikzset{b/.style={draw,rounded corners,align=center,text width=3.25cm,minimum height=1.05cm}}
\node[b] (id) at (0,0) {Identidad de ejecución,\\archivos y tamaños};
\node[b] (req) at (4.8,1.15) {Grafo y ubicación\\$M_{\mathrm{req}}$ por destino};
\node[b] (obs) at (4.8,-1.15) {Contratos y registros\\$M_{\mathrm{obs}}$ por entrega};
\node[b] (out) at (9.6,0) {Consistencia y cobertura\\DME y desglose};
\draw[->] (id)--(req);\draw[->] (id)--(obs);\draw[->] (req)--(out);\draw[->] (obs)--(out);
\end{tikzpicture}\end{lrbox}\fitdmediagram
\end{figure}

\begin{table}[htbp]\centering\small
\caption{Decisiones del modelo de medición}
\begin{tabularx}{\linewidth}{@{}P{2.65cm}YY@{}}
\toprule Decisión & Alternativa considerada & Razón de la decisión\\\midrule
Relevo obligatorio & Referencia sin relevo. & Omitir el coordinador subestima el volumen exigido cuando un archivo cruza entre workers.\\
Referencia por máquina & Contar una copia por consumidor. & Se cuenta una disponibilidad por worker y el relevo obligatorio en el coordinador; no una entrada por sandbox.\\
Ubicación efectiva & Optimizar sobre todas las asignaciones posibles. & La métrica describe la ejecución estudiada. La cota depende de sus ubicaciones, finales y del relevo por el coordinador.\\
Entrega científica completa & Usar el total bruto de transporte. & Los archivos del grafo y las entregas corroboradas se comparan en una misma unidad y frontera.\\
Razón con sus volúmenes & Clasificar configuraciones solo por DME. & El numerador y el denominador permiten identificar qué originó un cambio del cociente.\\
\bottomrule
\end{tabularx}
\end{table}

\chapter{Implementación de la solución}
\label{ch:implementacion}
\section{Componentes y tecnologías}
\Analyzer{} es una herramienta en Python que, después del workflow, combina el directorio de Pegasus con el historial exportado de HTCondor. Produce resúmenes y desgloses de obligaciones, entregas, evidencia y ubicación; cada análisis conserva la versión y la identidad del flujo.

El cálculo incorpora el relevo del capítulo~\ref{ch:diseno}: si un intermedio tiene consumidor en otro worker, añade la obligación al coordinador y la deduplica cuando este también es destino final. Si no se identifica el coordinador, marca incompleta la evidencia.

\paragraph{Integración de las fuentes.}
El grafo expresa tareas y archivos; el plan añade actividades de preparación y entrega y puede reunir tareas en un job. La lista secuencial, los comandos y las dependencias permiten reconstruir cada grupo; el historial asigna a sus tareas la máquina del job. La herramienta coteja flujo e historial, identifica el intento y comprueba terminación y eventos según el perfil.

Catálogos y metadatos permiten conciliar tamaños e identidades. Fuentes incompatibles, colisiones de rutas, ubicación ausente o tareas repetidas impiden reconstruir la parte afectada; nombres similares no bastan para establecer identidad.

\section{Conciliación de entregas}
Una dirección corresponde a la entrada de archivos al job o a su salida. Para cada job $j$ y dirección $d$, se reconstruye la lista efectiva de archivos científicos $S_{jd}$ y se separan los auxiliares. Si se conocen sus tamaños, el residuo de control se calcula restando del total nativo $b_{jd}$ los bytes auxiliares $a_{jd}$ y los tamaños de los archivos científicos:
\begin{equation}
r_{jd}=b_{jd}-a_{jd}-\sum_{f\in S_{jd}}s(f).
\label{eq:residuo}
\end{equation}
El perfil del bloque A exige residuo cero, conteos concordantes y consistencia con terminación e intento. Se contrastan los totales nativos de HTCondor por dirección con el último intento y se corrobora la lista mediante evidencia agregada del job. El caso GL3 del capítulo~\ref{ch:validacion} muestra la cuenta completa.

Si se desconoce algún tamaño auxiliar, el modo \texttt{SUCCESSFUL\_MANIFEST} exige manifiesto compatible, tamaños científicos y ejecución satisfactoria; el residuo queda indeterminado. La auditoría histórica del bloque A usa \texttt{EXACT\_BYTES}. Los archivos sin evidencia suficiente se conservan como faltantes para decidir el cálculo.

\paragraph{Reportes y control de consistencia.}
El resumen presenta $M_{\mathrm{req}}$, $M_{\mathrm{obs}}$, DME y estado; los desgloses identifican destinos, entregas por dirección y conciliación de auxiliares. Los diagnósticos separan evidencia faltante, contradicciones y perfiles fuera de alcance.

La comparación con código separado comprueba consistencia, aunque comparta definición y fuentes.

\paragraph{Trazabilidad de una ejecución agrupada.}
\label{sec:traza_impl}
El grupo se reconstruye antes de calcular ambos volúmenes. En GL3, plan e historial sitúan tres tareas en un job de W1, y el grafo conserva dos dependencias intermedias. La rama de obligaciones asigna origen y destino en W1 a esos archivos; la de entregas recoge los intermedios y el final declarados en el contrato de salida. Así, un intermedio puede no crear una obligación entre máquinas y aun aportar al observado mediante un retorno corroborado.



El siguiente capítulo verifica estas cuentas y examina qué entregas cambian al modificar agrupación o contratos.

\chapter{Validación de la solución}
\label{ch:validacion}
\section{Diseño de los contrastes}
La validación examina la correctitud de los volúmenes reconstruidos y la aplicación del mismo procedimiento en los cinco patrones. El diseño sigue el enfoque de \citet{jain1991} para definir la pregunta, la carga, los factores, las respuestas y los controles. La agrupación se estudia mediante tres mecanismos: eliminar reenvíos entre etapas, reutilizar un archivo entre consumidores y reunir tareas que consumen particiones distintas. Se examina cuándo DME aumenta, cuándo permanece constante y qué entregas explican el resultado. Los contrastes de tamaño y ubicación permiten precisar la interpretación de la referencia.

El bloque A histórico reúne 24 condiciones de la evaluación de DME con flujos sintéticos bajo CondorIO. Sus 72 ejecuciones elegibles corresponden a tres repeticiones por condición, enviadas en orden fijo. La unidad de observación es una ejecución completa. La selección posterior de quince condiciones y 45 ejecuciones organiza siete contrastes en los cinco patrones. AC1, AC2 y AC4 designan los niveles de uno, dos y cuatro consumidores por job en el caso de archivo compartido. El paquete histórico conserva sus identificadores originales; aquí se usan AC1, AC2 y AC4. R01-E y R02-E conservan sus identificadores históricos dentro del bloque A.

\subsection{Decisiones del diseño según Jain}
\label{sec:jain_diseno}
El enfoque de \citet{jain1991} y su material docente \citep{jain2017} orientan la definición del sistema, la carga, las respuestas, los factores y los controles. El servicio del sistema es ejecutar un grafo y entregar sus resultados declarados. Se usan cinco patrones con archivos de tamaños e identidades conocidos. En distribución se distinguen particiones diferentes y consumidores de un mismo archivo. Las respuestas principales son los volúmenes en bytes, DME y los conteos. La duración del flujo y los tiempos posteriores se miden por separado.

Cada contraste tiene una línea base y varía el tamaño, la agrupación o la ubicación según la pregunta. Se conserva el resto de los controles pertinentes y se verifica su valor efectivo con el plan y el historial. Los niveles de las fichas explican qué cambia y por qué. Los contrastes permiten comprobar el cálculo por archivo y destino. Un diseño factorial permitiría examinar interacciones entre los factores con niveles comparables en todos los patrones. El requerido se deriva de la definición; las cuentas del observado se reconstruyeron retrospectivamente desde contratos preservados y no constituyen predicciones previas al experimento.

\subsection*{Cobertura de la selección}
La Tabla~\ref{tab:fuera_contrastes} identifica las nueve condiciones elegibles que quedaron fuera de las fichas. La selección retrospectiva permite comprobar las cuentas de cada contraste. Estimar la frecuencia de estos comportamientos en otros flujos requeriría otro diseño. Las otras condiciones documentan asignaciones alternativas, cambios en el número de tareas, un nivel adicional de agrupación o cambios de pool y ubicación. Sus 27 ejecuciones se conservan como archivo de campaña y no integran las tablas de DME de esta tesis. El cotejo detallado de 348 direcciones de este capítulo corresponde a las 45 ejecuciones seleccionadas.

\begin{table}[htbp]\centering\small
\caption{Condiciones elegibles fuera de las fichas de validación}
\label{tab:fuera_contrastes}
\begin{tabularx}{\linewidth}{@{}P{3.3cm}Y@{}}
\toprule Condiciones & Qué añaden al bloque A\\\midrule
L02, D01 & Asignaciones alternativas del pipeline y de distribución.\\
L05 & Pipeline más largo, de cinco etapas.\\
A02B, D02B, R02B & Grafos de agregación, distribución y redistribución con menos ramas.\\
GD2 & Nivel intermedio entre D02 y GD4 para agrupar particiones distintas.\\
W01, W04 & Cambios conjuntos de asignación y número de máquinas disponibles.\\
\bottomrule
\end{tabularx}
\end{table}


\subsection{Líneas base y controles comunes}
La línea base es la configuración ejecutada con la cual se compara una variante. El volumen requerido se calcula por separado para cada ejecución. La sección~\ref{sec:repro_patrones} identifica las bases y variantes de los siete contrastes, con su justificación, cuenta esperada y decisión respaldada por los datos.



Los 10 MiB permiten seguir las entregas del pipeline. En procesamiento se duplica ese tamaño a 20 MiB. En I40 se cambia solo un intermedio a 40 MiB. Estos niveles sirven para comprobar cambios de volumen en cargas controladas. El tiempo se mide por separado.

En los contrastes históricos de agrupación se conservan el grafo, los tamaños, la asignación y la política de entrega. En los de ubicación se mantienen los archivos y la organización en jobs. R02-E modifica la asignación de varias etapas a la vez. Su efecto sobre el requerido se desglosa por dependencia, mientras que la duración corresponde a la configuración completa.

\subsection{Entorno y orden de toma de datos}
El bloque A se ejecutó con Pegasus 5.1.2, HTCondor 26.0.1 (compilación DAILY) y CondorIO en máquinas virtuales Ubuntu Server 24.04 LTS x64, en una misma región de la nube. Los manuales citados describen Pegasus 5.1.3. Los contratos se verificaron con los planes y registros de 5.1.2; no se presupone equivalencia entre versiones ni estabilidad fuera de la compilación ensayada. L01-r01 se excluyó por un cambio de intervalos durante la ejecución (sección~\ref{sec:protocolo_a}). El coordinador tiene dos CPU virtuales y cerca de 6,77 GiB visibles. Cada worker tiene una CPU virtual y cerca de 3,33 GiB visibles. El slot anuncia una CPU y 1410 MiB, y los jobs científicos solicitan una CPU y 256 MiB. Los contrastes seleccionados tienen dos workers disponibles y cada worker anuncia un slot. Se fijó la ubicación en el archivo de envío de cada job científico mediante \texttt{requirements = TARGET.Machine == "dme-worker1"} para W1 o el nombre \texttt{dme-worker2} para W2. Al agrupar, el requisito se aplica al job que reúne las tareas. Se contrastaron estas declaraciones con la máquina registrada en el historial. El orden de envío de los workflows y la ubicación de sus jobs son decisiones distintas. Cada ejecución se registra con su condición, repetición e identidad.



\subsection{Procedimiento reproducible de toma de datos}
\label{sec:protocolo_a}
El procedimiento mantiene unidos el caso previsto y lo que efectivamente sucedió. La revisión de cada ejecución sigue este orden.
\begin{enumerate}
\item Registrar condición, repetición, entradas, intermedios y finales. Identificar qué contraste usa el caso y qué controles deben coincidir con su base.
\item Revisar versiones, recursos, pool y cola. Corroborar qué máquinas están disponibles y qué requisitos se declararon. Anotar la posición del envío dentro de la secuencia.
\item Revisar el plan antes del envío. Comprobar qué tareas pertenecen a cada job y qué dependencias, archivos y finales declara. La agrupación se confirma en las descripciones efectivas.
\item Ejecutar y preservar el identificador, las marcas de inicio y fin de DAGMan, la ubicación e intento de cada job, y los contratos y registros de entrada y salida.
\item Reconstruir las dos cuentas. Para cada obligación, identificar el archivo, su tamaño y el destino. Para cada entrega, vincular archivo y dirección con el intento y conciliar conteos y bytes. Registrar los conflictos de cobertura o identidad.
\item Resolver elegibilidad y comparar con la base. Conservar por separado los casos que incumplen el perfil y sus diagnósticos. Presentar ambos volúmenes, la razón exacta y la duración del mismo workflow.
\end{enumerate}

La duración del workflow comienza y termina según las marcas de DAGMan. Las tareas de preparación y entrega que ocurren dentro de ese intervalo forman parte de la duración. La exportación del historial y la invocación de \Analyzer{} posteriores al término se cronometran por separado. Los registros de transferencias corroboran las entregas y sus volúmenes. Su resolución temporal no se utiliza aquí para estimar una tasa de red a partir de intervalos breves.

\begin{table}[htbp]\centering\small
\caption{Reglas de elegibilidad y comprobaciones del análisis}
\begin{tabularx}{\linewidth}{@{}P{3.3cm}YY@{}}
\toprule Comprobación & Criterio & Consecuencia si falla\\\midrule
Condición y controles & Coinciden grafo, tamaños, grupos, ubicación y política con el nivel documentado. & El caso queda fuera de la comparación del perfil fijo.\\
Terminación e intento & El workflow termina y el intento efectivo puede identificarse para cada job. & Se conserva el diagnóstico de ejecución y se revisa la elegibilidad.\\
Cobertura de archivos & Se reconstruyen los archivos necesarios y se respaldan las entregas con evidencia. & El volumen afectado y DME pueden quedar indeterminados.\\
Coincidencia del cálculo & Se comparan las cuentas manuales, el núcleo y el auditor local sobre los mismos insumos; se distingue este cotejo de la elegibilidad de la ejecución. & Se localiza la discrepancia antes de interpretar un contraste.\\
Uso de la duración & Se conserva el tiempo de cada ejecución y la organización efectiva de jobs. & Se acota la conclusión a la evidencia disponible.\\
\bottomrule
\end{tabularx}
\end{table}

De los 74 envíos del bloque A, dos quedaron fuera del perfil fijo. L01-r01 terminó y fue auditado, pero se ajustaron los intervalos del planificador y del negociador mientras el flujo seguía en curso. Sus repeticiones elegibles son r02, r03 y r04. L05-r01 se interrumpió al detener la sesión y quedó sin aceptación. Los dos registros se conservan. Las restantes 72 ejecuciones cumplieron las reglas de inclusión.


\section{Correctitud del resultado}
\subsection{Reconstrucción manual de la referencia}
Para comprobar que la agrupación conserva la referencia, se toma el pipeline de la Figura~\ref{fig:pat-pipeline}. La entrada $A$ reside en el coordinador. Las tres tareas están en W1. Los dos intermedios nacen en W1 y se consumen en W1. El único final, $B$, debe quedar en el coordinador. Todos los archivos miden 10 MiB. Por la ecuación~\eqref{eq:requerido}, la entrada y el final generan las únicas obligaciones por destino distinto de su origen. El requerido resulta de sumar $10+0+0+10=20$ MiB en las tres configuraciones.



\subsection{Comprobación del observado en las tres agrupaciones}
Los contratos efectivos corroborados permiten separar entradas y retornos, como muestra la Tabla~\ref{tab:direcciones_pipeline}. En el caso sin agrupar, cada intermedio se devuelve al coordinador y vuelve a entrar al siguiente job. El primer reenvío desaparece al reunir $T_1$ y $T_2$. Al agrupar las tres tareas desaparece también el segundo reenvío. Los retornos de los intermedios permanecen en los contratos observados. Este desglose explica por qué GL3 entrega 40 MiB, aun cuando su cota es de 20 MiB.

\begin{table}[htbp]\centering\small
\caption{Entregas corroboradas del pipeline según su agrupación}
\label{tab:direcciones_pipeline}
\begin{tabularx}{\linewidth}{@{}P{2.05cm}YYr@{}}
\toprule Caso & Entradas a W1 & Retornos al coordinador & MiB\\\midrule
L01-r02 & $A$ a $T_1$, $f_{12}$ a $T_2$ y $f_{23}$ a $T_3$. & $f_{12}$, $f_{23}$ y $B$. & 60\\
GL2-r01 & $A$ al grupo y $f_{23}$ a $T_3$. & $f_{12}$ y $f_{23}$ del grupo, y $B$ de $T_3$. & 50\\
GL3-r01 & $A$ al único grupo. & $f_{12}$, $f_{23}$ y $B$ del grupo. & 40\\
\bottomrule
\end{tabularx}
\end{table}

L01-r02 fue la primera repetición elegible de L01. La salida del job de GL3-r01 permite verificar un ejemplo del control de bytes. El contrato declara los tres archivos científicos, que aportan 31\,457\,280 bytes en conjunto. Los auxiliares aportan 14\,961 bytes y el contador nativo registra 31\,472\,241 bytes. Al sustituir en la ecuación~\eqref{eq:residuo}, el residuo es cero. El conteo nativo y el esperado coinciden en cinco archivos, tres científicos y dos auxiliares. La auditoría confirma además la terminación y el intento. Por tanto, esta dirección aporta 30 MiB al observado. La entrada científica aporta los otros 10 MiB y $\DME=20/40=1/2$.

\subsection{Comprobación del destino compartido}
\label{sec:validacion_compartido}
El caso AC1 comprueba la deduplicación por destino: $g$ requiere un retorno al coordinador y una entrega a W2, aunque allí lo consuman cuatro tareas. Su aporte es $10+10=20$ MiB; la entrada y los cuatro finales añaden 50 MiB. La cota de AC1, AC2 y AC4 es 70 MiB. Las entradas repetidas a sandboxes distintos se cuentan en el observado, no como nuevos destinos de la cota.

El desglose del auditor permite comprobar que las copias son entregas distintas y no registros duplicados. En AC1 aparecen cuatro entradas de $g$ en cuatro jobs de W2. En AC2 quedan dos entradas y en AC4 queda una. El retorno de $g$ desde el productor al coordinador sigue presente. En los tres casos, el plan y el historial sitúan a los cuatro consumidores en W2. La obligación por destino proviene de la definición. El desglose preservado distingue las copias; Analyzer y el auditor local calculan la cota con relevo.

\subsection{Cálculo manual y cotejo de las 45 ejecuciones}
Para las quince condiciones se reconstruyó el requerido por archivo y destino con el relevo obligatorio. En agregación, A02 suma $40+2(40)+40=160$ MiB; en particiones, D02 y GD4 suman $10+2(10)+10=40$ MiB; en archivo compartido, los tres niveles suman $10+2(10)+40=70$ MiB; en redistribución, R02-E suma $40+2(40+40+40)+40=320$ MiB. En los casos sin dependencia entre workers no se añade relevo. Se revisaron las quince condiciones y las nueve comparaciones libres más el piloto sin modificar el observado.

Las 348 direcciones del bloque A conservan su conciliación histórica de bytes y conteos. El cálculo con relevo coincide en 45 registros de configuraciones documentadas y diez desgloses del contraste libre. Además, se reextrajeron las fuentes administrativas de las 45 cápsulas seleccionadas: Analyzer y una reconstrucción separada coinciden en volúmenes, obligaciones y cobertura. Las copias utilizadas omiten contenido científico y paquetes auxiliares; esta reextracción corrobora por manifiesto exitoso y no repite la conciliación exacta de auxiliares ni verifica el contenido omitido. Materiales distingue las dos comprobaciones.


\begin{table}[htbp]\centering\small
\caption{Volumen requerido y observado por condición (MiB)}
\label{tab:referencia_relevo}
\begin{tabular}{@{}lrrr@{}}
\toprule Caso & $M_{\mathrm{req}}$ & $M_{\mathrm{obs}}$ & DME\\\midrule
P01 & 20 & 20 & $1$\\
P02 & 40 & 40 & $1$\\
L01 & 20 & 60 & $1/3$\\
GL2 & 20 & 50 & $2/5$\\
GL3 & 20 & 40 & $1/2$\\
I40 & 20 & 120 & $1/6$\\
A01 & 80 & 160 & $1/2$\\
A02 & 160 & 160 & $1$\\
D02 & 40 & 40 & $1$\\
GD4 & 40 & 40 & $1$\\
AC1 & 70 & 100 & $7/10$\\
AC2 & 70 & 80 & $7/8$\\
AC4 & 70 & 70 & $1$\\
R01-E & 80 & 320 & $1/4$\\
R02-E & 320 & 320 & $1$\\
\bottomrule\end{tabular}\end{table}

\subsection{Antecedente de Montage}
\label{sec:montage}
Un piloto y tres ejecuciones históricas situaron \texttt{mProject} en W1 y \texttt{mShrink} en W2. Entrada y cabecera suman 1\,650\,491 bytes, el intermedio 529\,920 y los finales 665\,280. El requerido con relevo es $1\,650\,491+2(529\,920)+665\,280=3\,375\,611$ bytes. Las cuatro ejecuciones corroboran seis obligaciones y seis entregas, con observado igual al requerido y DME igual a uno. Este recálculo histórico no compara agrupaciones; se separa del contraste nuevo con ambas tareas en W1.

\section{Comprobación prospectiva de agrupación}
\label{sec:prospectiva}
La muestra fijada antes de ejecutar comprende 63 ejecuciones principales en doce condiciones: seis por condición sintética y tres por variante de Montage. Se completó la matriz original sin añadir condiciones ni ampliar la muestra por sus resultados. Las hipótesis se informaron con evidencia histórica; no se trata de una evaluación ciega de flujos nuevos. Los doce pilotos válidos se analizan aparte. Un intento de piloto Montage falló antes de iniciar la aplicación por saltos CRLF en el lanzador; se preservó sin DME. La corrección de formato usó identificadores nuevos y mantuvo el trabajo científico. No hubo reposiciones automáticas ni exclusiones por contrariar una predicción.

Dentro de cada familia se conservan grafo, entradas, tamaños, programa, finales, recursos y ubicación. Se usa el entorno descrito para A: Pegasus 5.1.2, HTCondor 26.0.1 y CondorIO. Cada job científico solicita una CPU y 256 MiB. El orden programado recorre una vez los seis órdenes de las ternas sintéticas y tres veces cada orden de particiones, con semilla 2026110809 y rotación de familias. Montage tiene tres bloques exploratorios, sin balance de sus seis órdenes. Las tres sesiones y sus pausas conservan el orden relativo previsto y se documentan en el libro de envíos; no se interpretan como control de carga o caché. Los pilotos comprobaron contratos y equivalencia antes de abrir la muestra principal.

Las variantes GL3-F y M-F intervienen además el contrato de retorno: mantienen los comandos, las dependencias y los finales, pero omiten intermedios consumidos dentro del único grupo. Se preserva el cambio exacto del archivo de envío. Son intervenciones distintas de agrupar, no una propiedad automática de toda agrupación.

Un auditor separado reconstruye obligaciones y entregas desde los planes, historiales y contratos, sin completar datos a partir de las predicciones. Después se coteja con el cálculo implementado. Las 312 direcciones de la muestra principal se conciliaron exactamente en bytes y conteos, incluidos los auxiliares. Los finales sintéticos coinciden por identidad, tamaño y SHA256 entre variantes de cada familia. En Montage se compararon además dimensiones, máscaras de valores finitos y píxeles; la diferencia máxima fue cero. Los pilotos y los datos principales se preservan por separado.

\begin{table}[htbp]\centering\small
\caption{Contrastes sintéticos: volúmenes previstos y medidos coincidentes (MiB); conteos por ejecución}\label{tab:prospectiva}
\begin{tabular}{@{}lrrrrrr@{}}
\toprule Caso & $n$ & $M_{\mathrm{req}}$ & $M_{\mathrm{obs}}$ & DME & Oblig. & Entregas\\\midrule
L01 & 6 & 20 & 60 & $1/3$ & 2 & 6\\
GL2 & 6 & 20 & 50 & $2/5$ & 2 & 5\\
GL3 & 6 & 20 & 40 & $1/2$ & 2 & 4\\
GL3-F & 6 & 20 & 20 & $1$ & 2 & 2\\
AC1 & 6 & 70 & 100 & $7/10$ & 7 & 10\\
AC2 & 6 & 70 & 80 & $7/8$ & 7 & 8\\
AC4 & 6 & 70 & 70 & $1$ & 7 & 7\\
D02 & 6 & 40 & 40 & $1$ & 13 & 13\\
GD4 & 6 & 40 & 40 & $1$ & 13 & 13\\
\bottomrule\end{tabular}
\end{table}
En las seis repeticiones por condición sintética se obtuvieron los mismos volúmenes y conteos de la Tabla~\ref{tab:prospectiva}; es lo esperado bajo contratos fijos. L01--GL3 reduce un tercio del observado (60 a 40 MiB); AC1--AC4 reduce 30\,\% (100 a 70 MiB). D02--GD4 conserva los 40 MiB y las trece entregas: reunir particiones distintas no elimina ninguna. Esto comprueba tres comportamientos de DME mediante entregas identificadas, sin inferir utilidad de la agrupación solo por el número de jobs.

GL3-F conserva el final y elimina los dos retornos de intermedios que permanecían en GL3: observado de 40 a 20 MiB, DME de $1/2$ a $1$. Esa diferencia de 20 MiB era evitable bajo este contrato y grafo. La intervención prueba su alcanzabilidad concreta; no demuestra que la cota sea alcanzable para cualquier flujo. A requerido común DME conserva el orden inverso del observado. Su aporte es interpretar la proximidad a las obligaciones, como muestra GL3 frente a GD4 con iguales 40 MiB observados y razones distintas.

\subsection{Agrupación del subworkflow de Montage}
MontagePy 2.3.0 \citep{good2025} ejecuta las mismas dos operaciones de reproyección y reducción de resolución en W1. M-S las separa; M-G las reúne; M-F además omite el retorno del intermedio \texttt{project.fits}. La imagen y la cabecera suman 1\,650\,491 bytes, el intermedio mide 529\,920 y los finales suman 665\,280. Como productor y consumidor están en W1, el intermedio no añade obligación: requerido $1\,650\,491+665\,280=2\,315\,771$ bytes en las tres variantes. Esta ubicación difiere del contraste histórico W1/W2, cuya referencia incluye relevo y destino remoto.

\begin{table}[htbp]\centering\small
\caption{Montage en W1: volúmenes previstos y medidos coincidentes (bytes); conteos por ejecución}\label{tab:montage_agrupacion}
\begin{tabular}{@{}lrrrrrr@{}}
\toprule Caso & $n$ & $M_{\mathrm{req}}$ & $M_{\mathrm{obs}}$ & DME & Oblig. & Entregas\\\midrule
M-S & 3 & 2\,315\,771 & 3\,375\,611 & $2315771/3375611$ & 4 & 6\\
M-G & 3 & 2\,315\,771 & 2\,845\,691 & $2315771/2845691$ & 4 & 5\\
M-F & 3 & 2\,315\,771 & 2\,315\,771 & $1$ & 4 & 4\\
\bottomrule\end{tabular}
\end{table}
M-G elimina el reenvío del intermedio: 529\,920 bytes menos que M-S (15,70\,\%). M-F elimina además su retorno y entrega 1\,059\,840 bytes menos que M-S (31,40\,\%), alcanzando la referencia sin cambiar los finales. Las tres repeticiones por condición coinciden en volúmenes y entregas. El caso comprueba esos mecanismos con datos y operaciones científicas reales; sigue siendo un subworkflow pequeño de dos operaciones, no un mosaico completo ni evidencia general sobre aplicaciones científicas.

\subsection{Límites de la comprobación}
Las pausas entre sesiones y los tres bloques exploratorios de Montage limitan la interpretación temporal. La duración se conserva por ejecución como respuesta secundaria. Agrupar cambia también jobs y actividad administrativa; no se atribuye una diferencia temporal exclusivamente a bytes. No se controló la caché fría ni la carga de fondo. La ubicación está fijada y contrastada con el historial, en un único entorno CondorIO. La preservación posterior enlaza copias nuevas idénticas con recibos de rutas, hashes y fechas previas, sin alterar sus bytes; no aporta copias físicas independientes. Los pilotos técnicos, el fallo preservado y los límites de sesión están declarados. El costo de captura nativa de OE4 sigue sin aislarse.

\section{Contraste con ubicación libre}
\label{sec:ubicacion_libre}
Los contrastes históricos de agrupación del pipeline fijaron las tareas en W1. Para observar cómo cambia la referencia cuando el planificador elige entre dos workers disponibles, se ejecutó un piloto GL2 y nueve comparaciones: L01, GL2 y GL3, tres repeticiones por agrupación. Se usó el mismo flujo sintético con archivos de 10 MiB. El orden predefinido de envío fue L01--GL2--GL3, GL3--L01--GL2 y GL3--GL2--L01. La secuencia organiza las rondas; no asigna al azar las tareas a workers.

Pegasus envió los jobs a un único sitio HTCondor. HTCondor seleccionó un worker elegible según la disponibilidad observada; por ello no se interpreta el selector de Pegasus como una política entre workers. La ubicación efectiva de cada tarea se recuperó del historial de ejecución y se usó para recalcular el requerido por destino. Las comparaciones libres se analizan separadas de las 45 ejecuciones históricas, que tuvieron ubicaciones controladas.

El piloto GL2 comprobó el procedimiento antes de las nueve comparaciones. $T_1$ y $T_2$ se agruparon y $T_3$ quedó en un job separado. Las tres tareas se observaron en W1; $M_{\mathrm{req}}=20$ MiB, $M_{\mathrm{obs}}=50$ MiB, DME $=2/5$ y la duración del DAG fue 106,935 s. El piloto se reporta aparte de las comparaciones.

La Tabla~\ref{tab:libre} registra la ubicación observada, ambos volúmenes, la razón exacta y la duración de cada comparación.
\begin{table}[htbp]\centering\small
\caption{Resultados de las nueve comparaciones con ubicación libre (volúmenes en MiB)}
\label{tab:libre}
\begin{tabular}{@{}P{1.05cm}cP{2.4cm}rrrc@{}}
\toprule Caso & Ronda & $T_1/T_2/T_3$ & $M_{\mathrm{req}}$ & $M_{\mathrm{obs}}$ & DME & DAG (s)\\\midrule
L01 & 1 & W1/W1/W1 & 20 & 60 & $1/3$ & 120,937\\
GL2 & 1 & W1/W1/W2 & 40 & 50 & $4/5$ & 100,796\\
GL3 & 1 & W1/W1/W1 & 20 & 40 & $1/2$ & 80,662\\
GL3 & 2 & W1/W1/W1 & 20 & 40 & $1/2$ & 81,659\\
L01 & 2 & W1/W2/W2 & 40 & 60 & $2/3$ & 120,933\\
GL2 & 2 & W1/W1/W1 & 20 & 50 & $2/5$ & 90,784\\
GL3 & 3 & W1/W1/W1 & 20 & 40 & $1/2$ & 70,646\\
GL2 & 3 & W1/W1/W1 & 20 & 50 & $2/5$ & 100,797\\
L01 & 3 & W1/W1/W1 & 20 & 60 & $1/3$ & 125,941\\
\bottomrule
\end{tabular}
\end{table}

En las tres repeticiones, $M_{\mathrm{obs}}$ fue 60 MiB para L01, 50 MiB para GL2 y 40 MiB para GL3. Las medianas de duración fueron 120,937, 100,796 y 80,662 s, respectivamente; GL3 tuvo la menor duración en cada ronda. Son resultados descriptivos de este entorno: cambian la agrupación y algunas ubicaciones, por lo que no se atribuye la diferencia temporal únicamente al volumen.

Con ubicación libre, L01 de la ronda 2 requiere $10+2(10)+0+10=40$ MiB y tiene DME $=2/3$. GL2 de la ronda 1 requiere $10+0+2(10)+10=40$ MiB y tiene DME $=4/5$. El piloto y las otras siete comparaciones mantienen 20 MiB requeridos. La cuenta incluye el retorno obligatorio cuando una dependencia cruza entre workers. GL2 de la primera ronda tiene mayor DME que GL3 ($4/5$ frente a $1/2$), pero entrega más bytes (50 frente a 40 MiB): entre ubicaciones distintas, DME no ordena el observado.

Para escoger una agrupación cuyo objetivo sea reducir bytes, el observado responde directamente a esa pregunta; el requerido y la ubicación explican las obligaciones que acompañan el resultado. DME ordena el volumen observado cuando se compara un requerido positivo común. En el bloque A, agrupar etapas que comparten una dependencia interna y consumidores de un mismo archivo redujo entregas; agrupar consumidores de particiones distintas no lo hizo. Los retornos intermedios registrados limitaron el ahorro. Las nueve comparaciones libres extienden esta lectura bajo una política común, sin generalizarla a otros flujos o planificadores.

\section{Consistencia y comparación entre patrones}
\label{sec:repro_patrones}
En las quince condiciones históricas del bloque A, las tres repeticiones coinciden en ubicación, tamaños y observado. Las fichas y tablas siguientes corresponden únicamente a esas 45 ejecuciones. Con esos insumos, el requerido y DME también coinciden. Es una consistencia esperable bajo contratos fijos, no una prueba independiente de correctitud. Las fichas muestran cómo cambian las obligaciones y las entregas en cinco patrones; las duraciones se presentan por separado.

\subsection{Procesamiento}
\label{sec:prueba_p}
\textbf{Pregunta y propósito.} \textquestiondown{}Conserva DME su valor al crecer los archivos con iguales obligaciones y entregas? Una tarea permite cotejar las entregas de entrada y final.

\begin{table}[htbp]\centering\small
\caption{Configuraciones del contraste de procesamiento}
\begin{tabularx}{\linewidth}{@{}P{1.4cm}YY@{}}
\toprule Caso & Configuración & Por qué se incluye\\\midrule
P01 & Una tarea en W1, una entrada externa de 10 MiB y un final de 10 MiB en el coordinador. & Base de dos obligaciones, sin intermedios.\\
P02 & La entrada y el final miden 20 MiB. La tarea, el worker y el plan se conservan. & Duplica ambos tamaños para comprobar la escala por bytes, no por número de archivos.\\
\bottomrule
\end{tabularx}
\end{table}

\textbf{Factor y controles.} Varía el tamaño uniforme de entrada y final (10 o 20 MiB); se mantienen una tarea por job, W1, destino, recursos y política de entrega. W2 sigue disponible. El nivel de 10 MiB coincide con la referencia del pipeline.

\textbf{Comprobación esperada.} Entrada y final llevan requerido y observado de 20 a 40 MiB; se esperan dos entregas y DME igual a uno en ambos casos. Se cotejan identidad y tamaño; el tiempo de DAGMan se registra por separado.

\textbf{Resultado y decisión.} P01 y P02 corroboraron dos entregas; requerido y observado pasaron de 20 a 40 MiB y DME se mantuvo en 1. La razón se conserva, aunque los volúmenes absolutos distinguen las cargas.

\subsection{Pipeline}
\label{sec:prueba_l}
\textbf{Pregunta y propósito.} Con flujo y asignación fijos, \textquestiondown{}qué agrupación entrega menos volumen? El pipeline compara una o dos dependencias internas; I40 aumenta un intermedio con el grafo fijo.

\begin{table}[htbp]\centering\small
\caption{Configuraciones del contraste de pipeline}
\begin{tabularx}{\linewidth}{@{}P{1.35cm}YY@{}}
\toprule Caso & Configuración & Por qué se incluye\\\midrule
L01 & $T_1$, $T_2$ y $T_3$ en tres jobs separados en W1. Entrada, dos intermedios y final de 10 MiB. & Base sin agrupación; expone ambas fronteras entre jobs.\\
GL2 & $T_1$ y $T_2$ en un job y $T_3$ en otro. Todas las tareas en W1. & Hace interna la primera dependencia y conserva la segunda frontera.\\
GL3 & Las tres tareas se ejecutan secuencialmente en un job en W1. & Elimina ambas fronteras y permite examinar los retornos declarados.\\
I40 & Mismo plan que L01, con $f_{12}$ de 40 MiB y los otros tres archivos de 10 MiB. & Aumenta cuatro veces un intermedio, con tareas y entregas previstas constantes.\\
\bottomrule
\end{tabularx}
\end{table}

\textbf{Factor y controles.} La agrupación cambia los jobs científicos (tres, dos, uno) con tamaños, tareas, dependencias, W1 y entrega constantes. L01 e I40 comparan tamaño en tres jobs.

\textbf{Comprobación esperada.} Los cuatro casos requieren 20 MiB por entrada y final. Con los retornos, L01, GL2 y GL3 implican seis, cinco y cuatro entregas de 10 MiB. En I40, retorno y reenvío añaden $2(40-10)=60$ MiB frente a L01; se cotejan con el plan y los registros.


\noindent\textbf{Resultado y decisión del pipeline.} L01, GL2 y GL3 observaron 60, 50 y 40 MiB; con 20 MiB requeridos, DME fue $1/3$, $2/5$ y $1/2$. Frente a L01, GL2 redujo el observado 16,7\,\% y GL3 33,3\,\%; GL3 tuvo el menor observado de las tres agrupaciones históricas con la misma política de entrega. Los dos retornos declarados explican el exceso.

En I40, el primer intermedio creció de 10 a 40 MiB. Retorno y reenvío añadieron 30 MiB cada uno: el observado pasó de 60 a 120 MiB. Productor y consumidor permanecen en W1; requerido siguió en 20 MiB y DME bajó de $1/3$ a $1/6$. El efecto depende del archivo y sus entregas; una DME menor no indica un fallo del cálculo.

\subsection{Agregación}
\label{sec:prueba_a}
\textbf{Pregunta y propósito.} Si el agregador cambia de máquina con los mismos productores y archivos, \textquestiondown{}cuánto cambia la referencia y cuánto cambia el observado? Distingue nuevas obligaciones de transferencias del mecanismo.

\begin{table}[htbp]\centering\small
\caption{Configuraciones del contraste de agregación}
\begin{tabularx}{\linewidth}{@{}P{1.4cm}YY@{}}
\toprule Caso & Configuración & Por qué se incluye\\\midrule
A01 & Cuatro productores y un agregador en W1, en cinco jobs. Cuatro entradas y cuatro intermedios de 10 MiB, y un final de 40 MiB. & Coloca productores y consumidor en el mismo sitio. Permite separar los retornos del mecanismo de las obligaciones entre máquinas.\\
A02 & Los cuatro productores se mantienen en W1. El agregador pasa a W2. Se conservan los cinco jobs y los archivos. & Cambia solo la ubicación del agregador y muestra el requerido por los nuevos destinos.\\
\bottomrule
\end{tabularx}
\end{table}

\textbf{Factor y controles.} Se varía el worker del agregador (W1 o W2); se mantienen la asignación de productores, los jobs, las identidades de archivos y la entrega final al coordinador. El final mide 40 MiB, igual a la suma de los cuatro intermedios.

\textbf{Comprobación esperada.} Entradas y final suman 80 MiB en A01. En A02, los cuatro intermedios de 10 MiB retornan al coordinador y se entregan a W2, por lo que se añaden $4(10+10)=80$ MiB y la cota es 160 MiB. El observado es 160 MiB en ambos casos.

\textbf{Resultado y decisión.} A01 tiene una cota de 80 MiB y un observado de 160 MiB: DME es $1/2$. En A02 ambos son 160 MiB y DME es uno; los retornos intermedios son relevo obligatorio. La cota cambia con la ubicación; ambas configuraciones empatan en volumen observado.

\subsection{Distribución}
\label{sec:prueba_d}
\textbf{Pregunta y propósito.} \textquestiondown{}Puede la agrupación conservar las entregas cuando los consumidores necesitan archivos distintos? Se contrasta con el pipeline, que tiene dependencias internas.

\begin{table}[htbp]\centering\small
\caption{Configuraciones del contraste de distribución}
\begin{tabularx}{\linewidth}{@{}P{1.4cm}YY@{}}
\toprule Caso & Configuración & Por qué se incluye\\\midrule
D02 & Un divisor en W1 y cuatro consumidores en W2, en jobs separados. Entrada de 10 MiB y cuatro particiones y finales de 2,5 MiB cada uno. & Es una base con cuatro consumidores de archivos distintos y con una ubicación que exige entregar cada partición a W2.\\
GD4 & El divisor permanece separado en W1 y los cuatro consumidores se ejecutan en un job en W2. & Agrupa todos los consumidores; comprueba si conserva las cuatro entregas.\\
\bottomrule
\end{tabularx}
\end{table}

\textbf{Factor y controles.} Se varía el número de jobs de consumo (cuatro o uno), con ubicaciones, grafo y política de entrega fijos. Cada consumidor lee una partición y devuelve un final del mismo tamaño. Las cuatro particiones de 2,5 MiB suman 10 MiB.

\textbf{Comprobación esperada.} La cota suma 10 MiB de entrada, 10 MiB de retorno de particiones, 10 MiB de entrega a W2 y 10 MiB de finales: 40 MiB. El observado también suma 40 MiB; agrupar no elimina las particiones distintas ni su relevo obligatorio.

\textbf{Resultado y decisión.} D02 y GD4 conservaron 40 MiB requeridos y observados, con DME igual a uno. Las trece entregas científicas permanecieron al pasar de cinco jobs a dos; como el observado ya coincide con la cota, agrupar consumidores de particiones distintas no redujo bytes. La ficha siguiente examina consumidores de un mismo archivo.

\subsection{Distribución de un archivo compartido}
\label{sec:prueba_b}
\textbf{Pregunta y propósito.} Cuando varios consumidores requieren un mismo archivo, \textquestiondown{}qué cambia al agruparlos? Se compara la reutilización con D02 y GD4, donde cada consumidor necesita una partición distinta.

\begin{table}[htbp]\centering\small
\caption{Configuraciones de la distribución de un archivo compartido}
\begin{tabularx}{\linewidth}{@{}P{1.3cm}YY@{}}
\toprule Caso & Configuración & Por qué se incluye\\\midrule
AC1 & Cuatro consumidores en jobs separados en W2; productor en W1. & Base con cuatro entradas del mismo archivo a W2.\\
AC2 & Dos jobs de consumo, uno para $C_1$ y $C_2$ y otro para $C_3$ y $C_4$. & Comprueba una entrada por grupo de dos consumidores.\\
AC4 & Un job de consumo para las cuatro tareas. & Comprueba si un grupo requiere una entrada.\\
\bottomrule
\end{tabularx}
\end{table}

\textbf{Factor y controles.} Varía el número de consumidores por job: uno, dos o cuatro. El productor permanece en W1 y los consumidores en W2. La entrada, $g$ y cada final miden 10 MiB; cada consumidor produce un final distinto que se entrega al coordinador.

\textbf{Comprobación esperada.} La cota es 70 MiB: 10 MiB de entrada a W1, 20 MiB por el retorno y la entrega de $g$, y 40 MiB de finales. El observado suma 60 MiB fijos (entrada, retorno de $g$ y finales) y una entrada de $g$ por job de consumo: 100, 80 o 70 MiB en total.

\textbf{Resultado y decisión.} AC1, AC2 y AC4 observaron 100, 80 y 70 MiB; DME fue 7/10, 7/8 y 1. AC4 alcanza la cota con una copia de $g$ a W2 y reduce el observado 30\,\% frente a AC1. La agrupación permite reutilizar $g$ dentro de cada job. En D02 y GD4, cada consumidor necesita una partición distinta, por lo que agrupar no elimina esas cuatro entradas.

\subsection{Redistribución}
\label{sec:prueba_r}
\textbf{Pregunta y propósito.} \textquestiondown{}Cómo se interpreta DME cuando una nueva asignación introduce destinos en varias etapas del mismo flujo? El flujo combina producción, agregación, división y consumo.

\begin{table}[htbp]\centering\small
\caption{Configuraciones del contraste de redistribución}
\begin{tabularx}{\linewidth}{@{}P{1.4cm}YY@{}}
\toprule Caso & Configuración & Por qué se incluye\\\midrule
R01-E & Cuatro productores, un agregador, un divisor y cuatro consumidores, todos en W1. Cada tarea tiene su job. & Es una base con tres clases de intermedios y todas las dependencias científicas locales en W1.\\
R02-E & Productores y divisor en W1. Agregador y consumidores en W2. Se conservan los diez jobs. & Introduce destinos entre máquinas en las tres etapas intermedias. Comprueba la referencia ante nuevas obligaciones.\\
\bottomrule
\end{tabularx}
\end{table}

\textbf{Factor y controles.} Se compara un factor de asignación de etapas con dos niveles completos. Cambian conjuntamente varias ubicaciones en R02-E. Se mantienen el grafo, las entradas y finales de 10 MiB, los cuatro primeros intermedios y las cuatro particiones de 10 MiB, el agregado de 40 MiB y la política de entrega. El cambio conjunto sirve para contrastar la suma por destino. Aislar efectos temporales requeriría otros contrastes.

\textbf{Comprobación esperada.} Las entradas y finales aportan 80 MiB. En R01-E no hay relevo intermedio obligatorio entre workers. En R02-E, los tres conjuntos intermedios aportan 40 MiB cada uno y requieren retorno y entrega: la cota es $80+2(40+40+40)=320$ MiB. El observado es 320 MiB en ambas configuraciones.

\textbf{Resultado y decisión.} R01-E tiene 80 MiB requeridos, 320 MiB observados y DME $=1/4$: existe exceso sobre la cota. R02-E tiene 320 MiB en ambos volúmenes y DME igual a uno: los recorridos intermedios son relevo obligatorio. No hay ahorro entre configuraciones; cambia la referencia por ubicación.



\subsection{Conteos que sustentan la comparación}
\label{sec:conteos_a}
La Tabla~\ref{tab:conteos_a} presenta las obligaciones calculadas con relevo y los conteos de ejecución preservados. La columna Jobs cuenta los jobs científicos. Los valores corresponden a una ejecución y coinciden en las tres repeticiones de cada condición.

\begin{center}
\begin{minipage}{\linewidth}
\centering\small
\captionsetup{hypcap=false}
\captionof{table}{Conteos por workflow en las condiciones seleccionadas}
\label{tab:conteos_a}
\begin{tabular}{@{}lrrrr@{}}
\toprule Caso & Jobs & Obligaciones & Entregas & Direcciones conciliadas\\\midrule
P01 & 1 & 2 & 2 & 2\\
P02 & 1 & 2 & 2 & 2\\
L01 & 3 & 2 & 6 & 6\\
GL2 & 2 & 2 & 5 & 4\\
GL3 & 1 & 2 & 4 & 2\\
I40 & 3 & 2 & 6 & 6\\
A01 & 5 & 5 & 13 & 10\\
A02 & 5 & 13 & 13 & 10\\
D02 & 5 & 13 & 13 & 10\\
GD4 & 2 & 13 & 13 & 4\\
AC1 & 5 & 7 & 10 & 10\\
AC2 & 3 & 7 & 8 & 6\\
AC4 & 2 & 7 & 7 & 4\\
R01-E & 10 & 8 & 26 & 20\\
R02-E & 10 & 26 & 26 & 20\\
\bottomrule
\end{tabular}
\end{minipage}
\end{center}

En A02, D02, GD4, AC4 y R02-E coinciden obligaciones y entregas, además de los bytes. Bajo cobertura inyectiva por identidad y destino, la igualdad de cardinalidad hace biyectiva la correspondencia: no quedan entregas adicionales. La igualdad de conteos o bytes por sí sola no descarta sustituciones. En GL3, dos direcciones contienen cuatro entregas; en GD4, agrupar conserva trece entregas.
\begin{table}[htbp]\centering\small
\caption{Duración total de los flujos por condición}
\label{tab:resultados}
\begin{tabular}{@{}lrrl@{}}
\toprule Caso & Ejecuciones & Mediana (s) & Rango (s)\\\midrule
P01 & 3 & 71,623 & 70,631--80,914\\
P02 & 3 & 86,641 & 70,630--91,646\\
L01 & 3 & 120,925 & 120,922--130,947\\
GL2 & 3 & 101,803 & 96,793--151,830\\
GL3 & 3 & 81,666 & 70,658--91,678\\
I40 & 3 & 120,948 & 115,952--125,939\\
A01 & 3 & 111,048 & 110,982--116,111\\
A02 & 3 & 131,184 & 126,975--136,069\\
D02 & 3 & 130,967 & 111,162--141,075\\
GD4 & 3 & 100,836 & 100,821--106,826\\
AC1 & 3 & 137,194 & 121,952--141,118\\
AC2 & 3 & 126,833 & 105,933--131,957\\
AC4 & 3 & 105,815 & 100,814--115,847\\
R01-E & 3 & 155,763 & 140,613--156,397\\
R02-E & 3 & 215,692 & 205,923--215,854\\
\bottomrule
\end{tabular}
\par\smallskip Los tiempos corresponden al DAG completo. Los volúmenes y DME se interpretan en las fichas de cada patrón.
\end{table}

\subsection{Comparación de duración dentro del bloque A}
\label{sec:duracion_a}
La Tabla~\ref{tab:resultados} resume las duraciones completas de DAGMan por condición \citep{htcondortiming}. La Figura~\ref{fig:tiempos_pipeline} muestra las tres ejecuciones de cada agrupación del pipeline. Los envíos siguieron un orden fijo y las repeticiones de configuraciones distintas no forman pares temporales.

\begin{figure}[htbp]
\caption{Duraciones individuales y medianas del pipeline}
\label{fig:tiempos_pipeline}
\centering
\begin{tikzpicture}[x=0.094cm,y=1cm,font=\small]
\foreach \x in {70,90,110,130,150} {
\draw[gray!28] (\x,0.30)--(\x,3.45);
\node[below] at (\x,0.26) {\x};
}
\node at (110,-0.50) {Duración de DAGMan (s)};
\foreach \y/\name in {3/L01,2/GL2,1/GL3} {
\node[left,font=\small\bfseries] at (66,\y) {\name};
}
\foreach \x/\y in {130.947/2.87,120.922/3.00,120.925/3.13,151.830/1.87,101.803/2.00,96.793/2.13,70.658/0.87,91.678/1.00,81.666/1.13} {
\fill[blue!55!black] (\x,\y) circle (2.25pt);
}
\foreach \x/\y in {120.925/3,101.803/2,81.666/1} {
\draw[teal!65!black,very thick] (\x,\y-0.28)--(\x,\y+0.28);
}
\end{tikzpicture}
\par\smallskip\footnotesize Puntos, ejecuciones individuales. Trazo, mediana. La separación vertical dentro de una misma fila facilita distinguir puntos próximos.
\end{figure}

Las medianas de L01 y GL2 fueron 120,925 y 101,803 s. El máximo de GL2, 151,830 s, superó el de L01, 130,947 s. Las tres duraciones de GL3 quedaron por debajo del mínimo de L01.

Para el archivo compartido, las medianas de AC1, AC2 y AC4 fueron 137,194, 126,833 y 105,815 s, a la vez que bajaron los bytes entregados. Los rangos de AC1 y AC2 se solapan. D02 y GD4 conservaron el volumen y sus medianas fueron 130,967 y 100,836 s. Las diferencias de duración describen las configuraciones completas. La agrupación también cambia el número de jobs, por lo que el diseño no aísla el efecto de los bytes.

\subsection{Criterios para recomendar una configuración}
\label{sec:decisiones_a}
La recomendación depende del criterio del usuario. Para reducir bytes científicos se comparan observados bajo trabajo y fronteras equivalentes. Cuando el requerido positivo es común, un DME mayor corresponde a un volumen observado menor. Si la ubicación cambia, el valor de DME se interpreta con los dos volúmenes y el desglose por destino. En ninguna de estas comparaciones el valor de la razón reemplaza la duración registrada.



\subsection{Sobrecosto de instrumentación}
\label{sec:costo_a}
La medición utiliza los planes y registros nativos de Pegasus y HTCondor. La exportación del historial y el cálculo de DME ocurren después del workflow; sus tiempos no estiman el costo de generar registros durante la ejecución. En la campaña, las medianas de exportación y análisis fueron 0,128 y 0,214 s, respectivamente, con el ejecutable de entonces. No se suman como mediana conjunta ni se atribuyen al código actual. El costo de captura nativa no quedó aislado y limita el cumplimiento de OE4. Los registros individuales se conservan en los materiales.



\subsection{Alcance de la reproducción}
La coincidencia entre repeticiones demuestra consistencia bajo las condiciones ensayadas. Los materiales permiten recalcular los resúmenes del bloque A a partir de datos procesados y ejecutar \Analyzer{} sobre cinco casos históricos portables. Reauditar las 45 ejecuciones desde sus fuentes requiere las cápsulas administrativas. Repetir la medición temporal exige además los ejecutables, sus parámetros y una configuración comparable del entorno.

La validación histórica usa flujos sintéticos pequeños bajo CondorIO; la nueva campaña añade contrastes prospectivos de agrupación y retorno, incluidos dos operadores de Montage. El contraste libre cubre solo el pipeline; la campaña prospectiva mantiene ubicación fija. Otros tamaños o mecanismos requieren contrastes propios. El costo de exportación y análisis se delimita en la sección~\ref{sec:costo_a}.

\chapter{Conclusiones}
DME responde la pregunta de investigación con una referencia por archivo y destino, condicionada a la ubicación efectiva y al relevo obligatorio de CondorIO. La razón expresa proximidad volumétrica a esa cota; el desglose explica las copias adicionales. Con requerido positivo común, mayor DME equivale a menor observado.

La campaña prospectiva completa las 63 ejecuciones principales previstas en doce condiciones, separadas de doce pilotos válidos y un fallo técnico preservado. Seis repeticiones por condición sintética comprueban tres comportamientos: pipeline de 60 a 50 y 40 MiB con 20 MiB requeridos; archivo compartido de 100 a 80 y 70 MiB con 70 MiB requeridos; particiones distintas en 40 MiB con DME igual a uno. Los cambios se explican por reenvíos y entradas repetidas.

Agrupar las tres etapas conserva retornos intermedios. Omitirlos mediante una intervención separada reduce el observado de 40 a 20 MiB y lleva DME de $1/2$ a $1$, manteniendo el final. En Montage, tres repeticiones por variante en W1 requieren 2\,315\,771 bytes. Los observados al separar, agrupar y omitir además el retorno intermedio son, respectivamente, 3\,375\,611, 2\,845\,691 y 2\,315\,771 bytes. Los finales conservan dimensiones, máscaras y píxeles. Se demuestra exceso evitable en esos dos grafos; el contraejemplo analítico de jobs aislados sin caché impide generalizar esa conclusión.

GL3 y GD4 entregan 40 MiB con requeridos distintos. DME distingue su relación con las obligaciones sin establecer superioridad global. El contraste histórico libre muestra además que una razón mayor puede acompañar más bytes. Para elegir por volumen se compara el observado; para interpretarlo se conservan la referencia y su desglose.

Los tres primeros objetivos quedan respaldados por la caracterización, definición y cálculo implementado. OE4 comprueba correctitud, consistencia, sensibilidad y límites, pero permanece parcial en el costo de captura nativa. La evidencia comprende flujos pequeños, dos operaciones de Montage y un único entorno CondorIO. Las duraciones describen configuraciones completas y no aíslan el efecto temporal de los bytes.

\chapter{Trabajo futuro}
La validación prospectiva comprueba tres mecanismos de agrupación y dos intervenciones de retorno bajo condiciones acotadas. Queda estudiar otros grafos y restricciones de reutilización para determinar cuándo la cota es alcanzable. El contraejemplo analítico sugiere distinguir la referencia por máquina del mínimo factible bajo aislamiento de jobs. No basta extrapolar desde las configuraciones con DME igual a uno.

Otro paso es comprobar tamaños y cargas mayores, manteniendo explícita la frontera de archivos completos y separando el cambio de ubicación del de contratos. Una evaluación temporal requeriría controlar carga y cachés y distinguir el efecto de los jobs del volumen entregado. También queda aislar el costo de generar registros nativos, pendiente de OE4. Los tiempos posteriores de exportación y cálculo no lo miden.

\chapter{Materiales}
El paquete reúne fuentes, código, pruebas y datos. Desde su raíz, en Python 3.12:
\begin{enumerate}
\item \texttt{python}\ \path{analysis/rebuild_relay_dataset.py}: genera \path{data/derived/dme_relevo_55.csv}, con volúmenes, obligaciones y procedencia.
\item \texttt{python}\ \path{analysis/audit_relay.py}: coteja esos 55 registros con código separado y produce \path{data/derived/auditoria_relevo_55.csv}.
\end{enumerate}
Los 55 registros separan 45 de configuraciones documentadas y diez del contraste libre, incluido el piloto. El archivo A--H conserva resultados históricos; no sustituye las tablas con relevo. Las otras 27 ejecuciones de A no se presentan como recalculadas.

La reextracción posterior de las 45 cápsulas administrativas está documentada en \path{data/validation_native_20261107/}: conserva resultados por ejecución, cobertura y modos de corroboración. Las 348 direcciones se corroboran por manifiesto exitoso en las copias utilizadas, que omiten contenido científico y paquetes auxiliares. Ese análisis no repite la conciliación exacta histórica de auxiliares. El archivo completo de investigación sigue siendo necesario para verificar el contenido omitido.

\path{analysis/rebuild_montage_relay.py} recalcula Montage desde \path{data/validation_montage_20261108/}, sin ejecutar las operaciones astronómicas. El recibo identifica las cuatro cápsulas y las fuentes cotejadas.

El paquete incluye cinco casos históricos portables. El README identifica sus entradas, comandos y límites; su corroboración es por manifiesto exitoso.

\path{data/prospective_20261108/} conserva matriz, predicciones previas, libro de envíos, resultados por UUID, fallo técnico y equivalencia. \path{analysis/prospective_20261108/rebuild.py} recalcula los volúmenes y conteos desde los desgloses extraídos, sin imputar predicciones. Las cápsulas separadas preservan entradas, finales, planes, registros y auxiliares. Los paquetes estáticos de Pegasus se conservan una vez por hash, vinculados mediante los manifiestos. Reauditar fuentes y contenido requiere ambos conjuntos; el paquete compacto solo permite reproducir los datos procesados. Los recibos documentan compresión sin pérdida y enlaces posteriores de copias idénticas, con rutas y bytes conservados.

\path{tools/verify_snapshot.py} comprueba integridad, no correctitud científica. El depósito privado y la licencia siguen pendientes; esta revisión no se ha publicado en GitHub.

\clearpage
\begin{thebibliography}{99}
\bibitem[Bharathi et al.(2008)]{bharathi2008} Bharathi, S., Chervenak, A., Deelman, E., Mehta, G., Su, M.-H., \& Vahi, K. (2008). Characterization of scientific workflows. En \emph{Proceedings of the 2008 Third Workshop on Workflows in Support of Large-Scale Science} (pp. 1--10). IEEE. \url{https://doi.org/10.1109/WORKS.2008.4723958}
\bibitem[Chen et al.(2015)]{chen2015} Chen, W., Ferreira da Silva, R., Deelman, E., \& Sakellariou, R. (2015). Using imbalance metrics to optimize task clustering in scientific workflow executions. \emph{Future Generation Computer Systems, 46}, 69--84. \url{https://doi.org/10.1016/j.future.2014.09.014}
\bibitem[Deelman et al.(2015)]{deelman2015} Deelman, E., Vahi, K., Juve, G., Rynge, M., Callaghan, S., Maechling, P. J., Mayani, R., Chen, W., Ferreira da Silva, R., Livny, M., \& Wenger, K. (2015). Pegasus, a workflow management system for science automation. \emph{Future Generation Computer Systems, 46}, 17--35. \url{https://doi.org/10.1016/j.future.2014.10.008}
\bibitem[Fortmeier et al.(2013)]{fortmeier2013} Fortmeier, O., Bücker, H. M., Fagginger Auer, B. O., \& Bisseling, R. H. (2013). A new metric enabling an exact hypergraph model for the communication volume in distributed-memory parallel applications. \emph{Parallel Computing, 39}(8), 319--335. \url{https://doi.org/10.1016/j.parco.2013.05.003}
\bibitem[Fuhrer et al.(2018)]{fuhrer2018} Fuhrer, O., Chadha, T., Hoefler, T., Kwasniewski, G., Lapillonne, X., Leutwyler, D., Lüthi, D., Osuna, C., Schär, C., Schulthess, T. C., \& Vogt, H. (2018). Near-global climate simulation at 1 km resolution: Establishing a performance baseline on 4888 GPUs with COSMO 5.0. \emph{Geoscientific Model Development, 11}, 1665--1681. \url{https://doi.org/10.5194/gmd-11-1665-2018}
\bibitem[Good(2025)]{good2025} Good, J. (2025). \emph{MontagePy} (Versión 2.3.0) [Software]. PyPI. \url{https://pypi.org/project/MontagePy/2.3.0/}
\bibitem[HTCondor Team(s. f.-a)]{htcondordagman} HTCondor Team. (s. f.-a). \emph{DAGMan workflows} (Rama 26.0). Recuperado el 28 de octubre de 2026, de \url{https://htcondor.readthedocs.io/en/v26_0-branch/automated-workflows/index.html}
\bibitem[HTCondor Team(s. f.-b)]{htcondortiming} HTCondor Team. (s. f.-b). \emph{Informational files: Workflow metrics} (Manual, rama 26.0). Recuperado el 30 de octubre de 2026, de \url{https://htcondor.readthedocs.io/en/v26_0-branch/automated-workflows/dagman-information-files.html}
\bibitem[HTCondor Team(s. f.-c)]{htcondorintro} HTCondor Team. (s. f.-c). \emph{Introduction to HTCondor} (Rama 26.0). Recuperado el 28 de octubre de 2026, de \url{https://htcondor.readthedocs.io/en/v26_0-branch/users-manual/welcome-to-htcondor.html}
\bibitem[HTCondor Team(s. f.-d)]{htcondorordinary} HTCondor Team. (s. f.-d). \emph{Submitting jobs without a shared file system: HTCondor's file transfer mechanism} (Manual, rama 26.0). Recuperado el 29 de octubre de 2026, de \url{https://htcondor.readthedocs.io/en/v26_0-branch/users-manual/file-transfer.html}
\bibitem[Jain(1991)]{jain1991} Jain, R. (1991). \emph{The art of computer systems performance analysis: Techniques for experimental design, measurement, simulation, and modeling}. Wiley-Interscience.
\bibitem[Jain(2017)]{jain2017} Jain, R. (2017). \emph{Common mistakes and how to avoid them} [Diapositivas de clase]. Washington University in St. Louis. \url{https://www.cse.wustl.edu/~jain/cse567-17/ftp/k_02mst.pdf}
\bibitem[Papadimitriou et al.(2021)]{papadimitriou2021} Papadimitriou, G., Wang, C., Vahi, K., Ferreira da Silva, R., Mandal, A., Liu, Z., Mayani, R., Rynge, M., Kiran, M., Lynch, V. E., Kettimuthu, R., Deelman, E., Vetter, J. S., \& Foster, I. (2021). End-to-end online performance data capture and analysis for scientific workflows. \emph{Future Generation Computer Systems, 117}, 387--400. \url{https://doi.org/10.1016/j.future.2020.11.024}
\bibitem[Pegasus WMS(s. f.-a)]{pegasustransfers} Pegasus WMS. (s. f.-a). \emph{Data transfers} (Versión 5.1.3). Recuperado el 28 de octubre de 2026, de \url{https://pegasus.isi.edu/documentation/user-guide/data-transfers.html}
\bibitem[Pegasus WMS(s. f.-b)]{pegasusstats} Pegasus WMS. (s. f.-b). \emph{Monitoring, debugging and statistics} (Versión 5.1.3). Recuperado el 7 de noviembre de 2026, de \url{https://pegasus.isi.edu/documentation/user-guide/monitoring-debugging-stats.html}
\bibitem[Pegasus WMS(s. f.-c)]{pegasusprovenance} Pegasus WMS. (s. f.-c). \emph{Submit directory details: Kickstart record} (Versión 5.1.3). Recuperado el 28 de octubre de 2026, de \url{https://pegasus.isi.edu/documentation/user-guide/submit-directory.html}
\bibitem[Pietri \& Sakellariou(2018)]{pietri2018} Pietri, I., \& Sakellariou, R. (2018). Scheduling data-intensive scientific workflows with reduced communication. En \emph{Proceedings of the 30th International Conference on Scientific and Statistical Database Management}. ACM. \url{https://doi.org/10.1145/3221269.3221298}
\end{thebibliography}

\end{document}
